%% APÊNDICE — Caderno Complementar (gerado por expandir_volume.py)

\chapter{Caderno Complementar de Prática}

No 5º ano, o Caderno Complementar reforça cada folha Kumon com uma folha extra de caligrafia, análise silábica e escrita. Use como atividade de casa, reforço ou avaliação de apoio; cada folha é reproduzível (plancheta).

%% COMPLEMENTAR C-01: DÍGRAFO CH

\plancheta{Folha Complementar C-01 --- Reforço da Folha E-01 (Dígrafo CH)}

\begin{exercicio}[Folha Complementar C-01 --- Reforço da Folha E-01 (Dígrafo CH)]

\metadata{C-01}{Silábico-Alfabético}{EF02LP03, EF02LP04, EF02LP07}{D1}{EF02LP04}{CNE/CEB nº 2/2012}

\kumontiming{2}{90}

\noindent\textbf{Nome:} \underline{\hspace{3.2cm}} \quad \textbf{Data:} \underline{\hspace{1.6cm}} \quad \textbf{Nota:} \underline{\hspace{0.9cm}}/10


\begin{enumerate}

  \item \textbf{Copie as palavras em letra cursiva:}
  \begin{center}
  chave \quad chapéu \quad chuva \quad chão
  \end{center}
  \begin{center}
  \underline{\hspace{2.4cm}} \quad \underline{\hspace{2.4cm}} \quad \underline{\hspace{2.4cm}} \quad \underline{\hspace{2.4cm}}
  \end{center}

  \item \textbf{Escreva as sílabas das palavras:}
  \begin{center}
  CHAVE $=$ \underline{\hspace{1.5cm}} \underline{\hspace{1.5cm}} \quad CHAPÉU $=$ \underline{\hspace{1.5cm}} \underline{\hspace{1.5cm}} \quad CHUVA $=$ \underline{\hspace{1.5cm}} \underline{\hspace{1.5cm}}
  \end{center}

  \item \textbf{Complete a tabela de sílabas:}

\begin{center}
\resizebox{\textwidth}{!}{%
\begin{tabular}{ccccc}
  \sil{CH}{A} & \sil{CH}{E} & \sil{CH}{I} & \sil{CH}{O} & \sil{CH}{U} \\
  \end{tabular}}
\end{center}

  \item \textbf{Copie as frases em letra cursiva:}
  \begin{center}
  A chave abriu a porta.\\
  \underline{\hspace{12cm}}\\
  A chuva molhou o chão.\\
  \underline{\hspace{12cm}}
  \end{center}

  \item \textbf{Desenhe a frase:}
  \begin{center}
  \begin{tikzpicture}
    \draw[thick, dashed] (0,0) rectangle (12,4);
  \end{tikzpicture}
  \end{center}

\end{enumerate}

\end{exercicio}

%% COMPLEMENTAR C-02: DÍGRAFO LH

\plancheta{Folha Complementar C-02 --- Reforço da Folha E-02 (Dígrafo LH)}

\begin{exercicio}[Folha Complementar C-02 --- Reforço da Folha E-02 (Dígrafo LH)]

\metadata{C-02}{Silábico-Alfabético}{EF02LP03, EF02LP04, EF02LP07}{D1}{EF02LP04}{CNE/CEB nº 2/2012}

\kumontiming{2}{90}

\noindent\textbf{Nome:} \underline{\hspace{3.2cm}} \quad \textbf{Data:} \underline{\hspace{1.6cm}} \quad \textbf{Nota:} \underline{\hspace{0.9cm}}/10


\begin{enumerate}

  \item \textbf{Copie as palavras em letra cursiva:}
  \begin{center}
  ilha \quad coelho \quad orelha \quad milho
  \end{center}
  \begin{center}
  \underline{\hspace{2.4cm}} \quad \underline{\hspace{2.4cm}} \quad \underline{\hspace{2.4cm}} \quad \underline{\hspace{2.4cm}}
  \end{center}

  \item \textbf{Escreva as sílabas das palavras:}
  \begin{center}
  ILHA $=$ \underline{\hspace{1.5cm}} \underline{\hspace{1.5cm}} \quad COELHO $=$ \underline{\hspace{1.5cm}} \underline{\hspace{1.5cm}} \quad ORELHA $=$ \underline{\hspace{1.5cm}} \underline{\hspace{1.5cm}}
  \end{center}

  \item \textbf{Complete a tabela de sílabas:}

\begin{center}
\resizebox{\textwidth}{!}{%
\begin{tabular}{ccccc}
  \sil{LH}{A} & \sil{LH}{E} & \sil{LH}{I} & \sil{LH}{O} & \sil{LH}{U} \\
  \end{tabular}}
\end{center}

  \item \textbf{Copie as frases em letra cursiva:}
  \begin{center}
  O coelho comeu a folha.\\
  \underline{\hspace{12cm}}\\
  O filho mora na ilha.\\
  \underline{\hspace{12cm}}
  \end{center}

  \item \textbf{Desenhe a frase:}
  \begin{center}
  \begin{tikzpicture}
    \draw[thick, dashed] (0,0) rectangle (12,4);
  \end{tikzpicture}
  \end{center}

\end{enumerate}

\end{exercicio}

%% COMPLEMENTAR C-03: DÍGRAFO NH

\plancheta{Folha Complementar C-03 --- Reforço da Folha E-03 (Dígrafo NH)}

\begin{exercicio}[Folha Complementar C-03 --- Reforço da Folha E-03 (Dígrafo NH)]

\metadata{C-03}{Silábico-Alfabético}{EF02LP03, EF02LP04, EF02LP07}{D1}{EF02LP04}{CNE/CEB nº 2/2012}

\kumontiming{2}{90}

\noindent\textbf{Nome:} \underline{\hspace{3.2cm}} \quad \textbf{Data:} \underline{\hspace{1.6cm}} \quad \textbf{Nota:} \underline{\hspace{0.9cm}}/10


\begin{enumerate}

  \item \textbf{Copie as palavras em letra cursiva:}
  \begin{center}
  ninho \quad banho \quad sonho \quad pinho
  \end{center}
  \begin{center}
  \underline{\hspace{2.4cm}} \quad \underline{\hspace{2.4cm}} \quad \underline{\hspace{2.4cm}} \quad \underline{\hspace{2.4cm}}
  \end{center}

  \item \textbf{Escreva as sílabas das palavras:}
  \begin{center}
  NINHO $=$ \underline{\hspace{1.5cm}} \underline{\hspace{1.5cm}} \quad BANHO $=$ \underline{\hspace{1.5cm}} \underline{\hspace{1.5cm}} \quad SONHO $=$ \underline{\hspace{1.5cm}} \underline{\hspace{1.5cm}}
  \end{center}

  \item \textbf{Complete a tabela de sílabas:}

\begin{center}
\resizebox{\textwidth}{!}{%
\begin{tabular}{ccccc}
  \sil{NH}{A} & \sil{NH}{E} & \sil{NH}{I} & \sil{NH}{O} & \sil{NH}{U} \\
  \end{tabular}}
\end{center}

  \item \textbf{Copie as frases em letra cursiva:}
  \begin{center}
  O ninho tem três ovos.\\
  \underline{\hspace{12cm}}\\
  Tomo banho de manhã.\\
  \underline{\hspace{12cm}}
  \end{center}

  \item \textbf{Desenhe a frase:}
  \begin{center}
  \begin{tikzpicture}
    \draw[thick, dashed] (0,0) rectangle (12,4);
  \end{tikzpicture}
  \end{center}

\end{enumerate}

\end{exercicio}

%% COMPLEMENTAR C-04: QUE E QUI

\plancheta{Folha Complementar C-04 --- Reforço da Folha E-04 (QUE e QUI)}

\begin{exercicio}[Folha Complementar C-04 --- Reforço da Folha E-04 (QUE e QUI)]

\metadata{C-04}{Silábico-Alfabético}{EF02LP03, EF02LP04, EF02LP07}{D1}{EF02LP04}{CNE/CEB nº 2/2012}

\kumontiming{2}{90}

\noindent\textbf{Nome:} \underline{\hspace{3.2cm}} \quad \textbf{Data:} \underline{\hspace{1.6cm}} \quad \textbf{Nota:} \underline{\hspace{0.9cm}}/10


\begin{enumerate}

  \item \textbf{Copie as palavras em letra cursiva:}
  \begin{center}
  queijo \quad quilo \quad quente \quad quero
  \end{center}
  \begin{center}
  \underline{\hspace{2.4cm}} \quad \underline{\hspace{2.4cm}} \quad \underline{\hspace{2.4cm}} \quad \underline{\hspace{2.4cm}}
  \end{center}

  \item \textbf{Escreva as sílabas das palavras:}
  \begin{center}
  QUEIJO $=$ \underline{\hspace{1.5cm}} \underline{\hspace{1.5cm}} \quad QUILO $=$ \underline{\hspace{1.5cm}} \underline{\hspace{1.5cm}} \quad QUENTE $=$ \underline{\hspace{1.5cm}} \underline{\hspace{1.5cm}}
  \end{center}

  \item \textbf{Leia e escreva o número de sílabas de cada palavra:}
  \begin{center}
  QUEIJO $\rightarrow$ \underline{\hspace{0.9cm}} \quad QUILO $\rightarrow$ \underline{\hspace{0.9cm}} \quad QUENTE $\rightarrow$ \underline{\hspace{0.9cm}} \quad QUERO $\rightarrow$ \underline{\hspace{0.9cm}}
  \end{center}

  \item \textbf{Copie as frases em letra cursiva:}
  \begin{center}
  O queijo está na mesa.\\
  \underline{\hspace{12cm}}\\
  Eu quero um quilo de pão.\\
  \underline{\hspace{12cm}}
  \end{center}

  \item \textbf{Desenhe a frase:}
  \begin{center}
  \begin{tikzpicture}
    \draw[thick, dashed] (0,0) rectangle (12,4);
  \end{tikzpicture}
  \end{center}

\end{enumerate}

\end{exercicio}

%% COMPLEMENTAR C-05: GUE E GUI

\plancheta{Folha Complementar C-05 --- Reforço da Folha E-05 (GUE e GUI)}

\begin{exercicio}[Folha Complementar C-05 --- Reforço da Folha E-05 (GUE e GUI)]

\metadata{C-05}{Silábico-Alfabético}{EF02LP03, EF02LP04, EF02LP07}{D1}{EF02LP04}{CNE/CEB nº 2/2012}

\kumontiming{2}{90}

\noindent\textbf{Nome:} \underline{\hspace{3.2cm}} \quad \textbf{Data:} \underline{\hspace{1.6cm}} \quad \textbf{Nota:} \underline{\hspace{0.9cm}}/10


\begin{enumerate}

  \item \textbf{Copie as palavras em letra cursiva:}
  \begin{center}
  guerra \quad guia \quad guitarra \quad água
  \end{center}
  \begin{center}
  \underline{\hspace{2.4cm}} \quad \underline{\hspace{2.4cm}} \quad \underline{\hspace{2.4cm}} \quad \underline{\hspace{2.4cm}}
  \end{center}

  \item \textbf{Escreva as sílabas das palavras:}
  \begin{center}
  GUERRA $=$ \underline{\hspace{1.5cm}} \underline{\hspace{1.5cm}} \quad GUIA $=$ \underline{\hspace{1.5cm}} \underline{\hspace{1.5cm}} \quad GUITARRA $=$ \underline{\hspace{1.5cm}} \underline{\hspace{1.5cm}}
  \end{center}

  \item \textbf{Leia e escreva o número de sílabas de cada palavra:}
  \begin{center}
  GUERRA $\rightarrow$ \underline{\hspace{0.9cm}} \quad GUIA $\rightarrow$ \underline{\hspace{0.9cm}} \quad GUITARRA $\rightarrow$ \underline{\hspace{0.9cm}} \quad ÁGUA $\rightarrow$ \underline{\hspace{0.9cm}}
  \end{center}

  \item \textbf{Copie as frases em letra cursiva:}
  \begin{center}
  A água está quente.\\
  \underline{\hspace{12cm}}\\
  A guitarra toca na festa.\\
  \underline{\hspace{12cm}}
  \end{center}

  \item \textbf{Desenhe a frase:}
  \begin{center}
  \begin{tikzpicture}
    \draw[thick, dashed] (0,0) rectangle (12,4);
  \end{tikzpicture}
  \end{center}

\end{enumerate}

\end{exercicio}

%% COMPLEMENTAR C-06: ENCONTRO BR

\plancheta{Folha Complementar C-06 --- Reforço da Folha E-06 (Encontro BR)}

\begin{exercicio}[Folha Complementar C-06 --- Reforço da Folha E-06 (Encontro BR)]

\metadata{C-06}{Silábico-Alfabético}{EF02LP03, EF02LP04, EF02LP07}{D1}{EF02LP04}{CNE/CEB nº 2/2012}

\kumontiming{2}{90}

\noindent\textbf{Nome:} \underline{\hspace{3.2cm}} \quad \textbf{Data:} \underline{\hspace{1.6cm}} \quad \textbf{Nota:} \underline{\hspace{0.9cm}}/10


\begin{enumerate}

  \item \textbf{Copie as palavras em letra cursiva:}
  \begin{center}
  braço \quad briga \quad brilho \quad brincar
  \end{center}
  \begin{center}
  \underline{\hspace{2.4cm}} \quad \underline{\hspace{2.4cm}} \quad \underline{\hspace{2.4cm}} \quad \underline{\hspace{2.4cm}}
  \end{center}

  \item \textbf{Escreva as sílabas das palavras:}
  \begin{center}
  BRAÇO $=$ \underline{\hspace{1.5cm}} \underline{\hspace{1.5cm}} \quad BRIGA $=$ \underline{\hspace{1.5cm}} \underline{\hspace{1.5cm}} \quad BRILHO $=$ \underline{\hspace{1.5cm}} \underline{\hspace{1.5cm}}
  \end{center}

  \item \textbf{Leia e escreva o número de sílabas de cada palavra:}
  \begin{center}
  BRAÇO $\rightarrow$ \underline{\hspace{0.9cm}} \quad BRIGA $\rightarrow$ \underline{\hspace{0.9cm}} \quad BRILHO $\rightarrow$ \underline{\hspace{0.9cm}} \quad BRINCAR $\rightarrow$ \underline{\hspace{0.9cm}}
  \end{center}

  \item \textbf{Copie as frases em letra cursiva:}
  \begin{center}
  O braço dói hoje.\\
  \underline{\hspace{12cm}}\\
  O livro está aberto.\\
  \underline{\hspace{12cm}}
  \end{center}

  \item \textbf{Desenhe a frase:}
  \begin{center}
  \begin{tikzpicture}
    \draw[thick, dashed] (0,0) rectangle (12,4);
  \end{tikzpicture}
  \end{center}

\end{enumerate}

\end{exercicio}

%% COMPLEMENTAR C-07: ENCONTRO CR

\plancheta{Folha Complementar C-07 --- Reforço da Folha E-07 (Encontro CR)}

\begin{exercicio}[Folha Complementar C-07 --- Reforço da Folha E-07 (Encontro CR)]

\metadata{C-07}{Silábico-Alfabético}{EF02LP03, EF02LP04, EF02LP07}{D1}{EF02LP04}{CNE/CEB nº 2/2012}

\kumontiming{2}{90}

\noindent\textbf{Nome:} \underline{\hspace{3.2cm}} \quad \textbf{Data:} \underline{\hspace{1.6cm}} \quad \textbf{Nota:} \underline{\hspace{0.9cm}}/10


\begin{enumerate}

  \item \textbf{Copie as palavras em letra cursiva:}
  \begin{center}
  cravo \quad crédito \quad criar \quad criança
  \end{center}
  \begin{center}
  \underline{\hspace{2.4cm}} \quad \underline{\hspace{2.4cm}} \quad \underline{\hspace{2.4cm}} \quad \underline{\hspace{2.4cm}}
  \end{center}

  \item \textbf{Escreva as sílabas das palavras:}
  \begin{center}
  CRAVO $=$ \underline{\hspace{1.5cm}} \underline{\hspace{1.5cm}} \quad CRÉDITO $=$ \underline{\hspace{1.5cm}} \underline{\hspace{1.5cm}} \quad CRIAR $=$ \underline{\hspace{1.5cm}} \underline{\hspace{1.5cm}}
  \end{center}

  \item \textbf{Leia e escreva o número de sílabas de cada palavra:}
  \begin{center}
  CRAVO $\rightarrow$ \underline{\hspace{0.9cm}} \quad CRÉDITO $\rightarrow$ \underline{\hspace{0.9cm}} \quad CRIAR $\rightarrow$ \underline{\hspace{0.9cm}} \quad CRIANÇA $\rightarrow$ \underline{\hspace{0.9cm}}
  \end{center}

  \item \textbf{Copie as frases em letra cursiva:}
  \begin{center}
  O cravo é uma flor.\\
  \underline{\hspace{12cm}}\\
  A criança gosta de criar.\\
  \underline{\hspace{12cm}}
  \end{center}

  \item \textbf{Desenhe a frase:}
  \begin{center}
  \begin{tikzpicture}
    \draw[thick, dashed] (0,0) rectangle (12,4);
  \end{tikzpicture}
  \end{center}

\end{enumerate}

\end{exercicio}

%% COMPLEMENTAR C-08: ENCONTRO DR

\plancheta{Folha Complementar C-08 --- Reforço da Folha E-08 (Encontro DR)}

\begin{exercicio}[Folha Complementar C-08 --- Reforço da Folha E-08 (Encontro DR)]

\metadata{C-08}{Silábico-Alfabético}{EF02LP03, EF02LP04, EF02LP07}{D1}{EF02LP04}{CNE/CEB nº 2/2012}

\kumontiming{2}{90}

\noindent\textbf{Nome:} \underline{\hspace{3.2cm}} \quad \textbf{Data:} \underline{\hspace{1.6cm}} \quad \textbf{Nota:} \underline{\hspace{0.9cm}}/10


\begin{enumerate}

  \item \textbf{Copie as palavras em letra cursiva:}
  \begin{center}
  dragão \quad madre \quad padrão \quad drama
  \end{center}
  \begin{center}
  \underline{\hspace{2.4cm}} \quad \underline{\hspace{2.4cm}} \quad \underline{\hspace{2.4cm}} \quad \underline{\hspace{2.4cm}}
  \end{center}

  \item \textbf{Escreva as sílabas das palavras:}
  \begin{center}
  DRAGÃO $=$ \underline{\hspace{1.5cm}} \underline{\hspace{1.5cm}} \quad MADRE $=$ \underline{\hspace{1.5cm}} \underline{\hspace{1.5cm}} \quad PADRÃO $=$ \underline{\hspace{1.5cm}} \underline{\hspace{1.5cm}}
  \end{center}

  \item \textbf{Leia e escreva o número de sílabas de cada palavra:}
  \begin{center}
  DRAGÃO $\rightarrow$ \underline{\hspace{0.9cm}} \quad MADRE $\rightarrow$ \underline{\hspace{0.9cm}} \quad PADRÃO $\rightarrow$ \underline{\hspace{0.9cm}} \quad DRAMA $\rightarrow$ \underline{\hspace{0.9cm}}
  \end{center}

  \item \textbf{Copie as frases em letra cursiva:}
  \begin{center}
  O dragão cuspiu fogo.\\
  \underline{\hspace{12cm}}\\
  A pedra é pesada.\\
  \underline{\hspace{12cm}}
  \end{center}

  \item \textbf{Desenhe a frase:}
  \begin{center}
  \begin{tikzpicture}
    \draw[thick, dashed] (0,0) rectangle (12,4);
  \end{tikzpicture}
  \end{center}

\end{enumerate}

\end{exercicio}

%% COMPLEMENTAR C-09: ENCONTRO FR

\plancheta{Folha Complementar C-09 --- Reforço da Folha E-09 (Encontro FR)}

\begin{exercicio}[Folha Complementar C-09 --- Reforço da Folha E-09 (Encontro FR)]

\metadata{C-09}{Silábico-Alfabético}{EF02LP03, EF02LP04, EF02LP07}{D1}{EF02LP04}{CNE/CEB nº 2/2012}

\kumontiming{2}{90}

\noindent\textbf{Nome:} \underline{\hspace{3.2cm}} \quad \textbf{Data:} \underline{\hspace{1.6cm}} \quad \textbf{Nota:} \underline{\hspace{0.9cm}}/10


\begin{enumerate}

  \item \textbf{Copie as palavras em letra cursiva:}
  \begin{center}
  fruta \quad frio \quad frente \quad fritar
  \end{center}
  \begin{center}
  \underline{\hspace{2.4cm}} \quad \underline{\hspace{2.4cm}} \quad \underline{\hspace{2.4cm}} \quad \underline{\hspace{2.4cm}}
  \end{center}

  \item \textbf{Escreva as sílabas das palavras:}
  \begin{center}
  FRUTA $=$ \underline{\hspace{1.5cm}} \underline{\hspace{1.5cm}} \quad FRIO $=$ \underline{\hspace{1.5cm}} \underline{\hspace{1.5cm}} \quad FRENTE $=$ \underline{\hspace{1.5cm}} \underline{\hspace{1.5cm}}
  \end{center}

  \item \textbf{Leia e escreva o número de sílabas de cada palavra:}
  \begin{center}
  FRUTA $\rightarrow$ \underline{\hspace{0.9cm}} \quad FRIO $\rightarrow$ \underline{\hspace{0.9cm}} \quad FRENTE $\rightarrow$ \underline{\hspace{0.9cm}} \quad FRITAR $\rightarrow$ \underline{\hspace{0.9cm}}
  \end{center}

  \item \textbf{Copie as frases em letra cursiva:}
  \begin{center}
  A fruta está madura.\\
  \underline{\hspace{12cm}}\\
  Hoje está frio.\\
  \underline{\hspace{12cm}}
  \end{center}

  \item \textbf{Desenhe a frase:}
  \begin{center}
  \begin{tikzpicture}
    \draw[thick, dashed] (0,0) rectangle (12,4);
  \end{tikzpicture}
  \end{center}

\end{enumerate}

\end{exercicio}

%% COMPLEMENTAR C-10: ENCONTRO GR

\plancheta{Folha Complementar C-10 --- Reforço da Folha E-10 (Encontro GR)}

\begin{exercicio}[Folha Complementar C-10 --- Reforço da Folha E-10 (Encontro GR)]

\metadata{C-10}{Silábico-Alfabético}{EF02LP03, EF02LP04, EF02LP07}{D1}{EF02LP04}{CNE/CEB nº 2/2012}

\kumontiming{2}{90}

\noindent\textbf{Nome:} \underline{\hspace{3.2cm}} \quad \textbf{Data:} \underline{\hspace{1.6cm}} \quad \textbf{Nota:} \underline{\hspace{0.9cm}}/10


\begin{enumerate}

  \item \textbf{Copie as palavras em letra cursiva:}
  \begin{center}
  grade \quad grande \quad grama \quad grato
  \end{center}
  \begin{center}
  \underline{\hspace{2.4cm}} \quad \underline{\hspace{2.4cm}} \quad \underline{\hspace{2.4cm}} \quad \underline{\hspace{2.4cm}}
  \end{center}

  \item \textbf{Escreva as sílabas das palavras:}
  \begin{center}
  GRADE $=$ \underline{\hspace{1.5cm}} \underline{\hspace{1.5cm}} \quad GRANDE $=$ \underline{\hspace{1.5cm}} \underline{\hspace{1.5cm}} \quad GRAMA $=$ \underline{\hspace{1.5cm}} \underline{\hspace{1.5cm}}
  \end{center}

  \item \textbf{Leia e escreva o número de sílabas de cada palavra:}
  \begin{center}
  GRADE $\rightarrow$ \underline{\hspace{0.9cm}} \quad GRANDE $\rightarrow$ \underline{\hspace{0.9cm}} \quad GRAMA $\rightarrow$ \underline{\hspace{0.9cm}} \quad GRATO $\rightarrow$ \underline{\hspace{0.9cm}}
  \end{center}

  \item \textbf{Copie as frases em letra cursiva:}
  \begin{center}
  A grade protege o jardim.\\
  \underline{\hspace{12cm}}\\
  O elefante é grande.\\
  \underline{\hspace{12cm}}
  \end{center}

  \item \textbf{Desenhe a frase:}
  \begin{center}
  \begin{tikzpicture}
    \draw[thick, dashed] (0,0) rectangle (12,4);
  \end{tikzpicture}
  \end{center}

\end{enumerate}

\end{exercicio}

%% COMPLEMENTAR C-11: ENCONTRO PR

\plancheta{Folha Complementar C-11 --- Reforço da Folha E-11 (Encontro PR)}

\begin{exercicio}[Folha Complementar C-11 --- Reforço da Folha E-11 (Encontro PR)]

\metadata{C-11}{Silábico-Alfabético}{EF02LP03, EF02LP04, EF02LP07}{D1}{EF02LP04}{CNE/CEB nº 2/2012}

\kumontiming{2}{90}

\noindent\textbf{Nome:} \underline{\hspace{3.2cm}} \quad \textbf{Data:} \underline{\hspace{1.6cm}} \quad \textbf{Nota:} \underline{\hspace{0.9cm}}/10


\begin{enumerate}

  \item \textbf{Copie as palavras em letra cursiva:}
  \begin{center}
  prato \quad primo \quad praça \quad prêmio
  \end{center}
  \begin{center}
  \underline{\hspace{2.4cm}} \quad \underline{\hspace{2.4cm}} \quad \underline{\hspace{2.4cm}} \quad \underline{\hspace{2.4cm}}
  \end{center}

  \item \textbf{Escreva as sílabas das palavras:}
  \begin{center}
  PRATO $=$ \underline{\hspace{1.5cm}} \underline{\hspace{1.5cm}} \quad PRIMO $=$ \underline{\hspace{1.5cm}} \underline{\hspace{1.5cm}} \quad PRAÇA $=$ \underline{\hspace{1.5cm}} \underline{\hspace{1.5cm}}
  \end{center}

  \item \textbf{Leia e escreva o número de sílabas de cada palavra:}
  \begin{center}
  PRATO $\rightarrow$ \underline{\hspace{0.9cm}} \quad PRIMO $\rightarrow$ \underline{\hspace{0.9cm}} \quad PRAÇA $\rightarrow$ \underline{\hspace{0.9cm}} \quad PRÊMIO $\rightarrow$ \underline{\hspace{0.9cm}}
  \end{center}

  \item \textbf{Copie as frases em letra cursiva:}
  \begin{center}
  O prato tem fruta.\\
  \underline{\hspace{12cm}}\\
  O primo mora longe.\\
  \underline{\hspace{12cm}}
  \end{center}

  \item \textbf{Desenhe a frase:}
  \begin{center}
  \begin{tikzpicture}
    \draw[thick, dashed] (0,0) rectangle (12,4);
  \end{tikzpicture}
  \end{center}

\end{enumerate}

\end{exercicio}

%% COMPLEMENTAR C-12: ENCONTRO TR

\plancheta{Folha Complementar C-12 --- Reforço da Folha E-12 (Encontro TR)}

\begin{exercicio}[Folha Complementar C-12 --- Reforço da Folha E-12 (Encontro TR)]

\metadata{C-12}{Silábico-Alfabético}{EF02LP03, EF02LP04, EF02LP07}{D1}{EF02LP04}{CNE/CEB nº 2/2012}

\kumontiming{2}{90}

\noindent\textbf{Nome:} \underline{\hspace{3.2cm}} \quad \textbf{Data:} \underline{\hspace{1.6cm}} \quad \textbf{Nota:} \underline{\hspace{0.9cm}}/10


\begin{enumerate}

  \item \textbf{Copie as palavras em letra cursiva:}
  \begin{center}
  trem \quad trator \quad triste \quad troca
  \end{center}
  \begin{center}
  \underline{\hspace{2.4cm}} \quad \underline{\hspace{2.4cm}} \quad \underline{\hspace{2.4cm}} \quad \underline{\hspace{2.4cm}}
  \end{center}

  \item \textbf{Escreva as sílabas das palavras:}
  \begin{center}
  TREM $=$ \underline{\hspace{1.5cm}} \underline{\hspace{1.5cm}} \quad TRATOR $=$ \underline{\hspace{1.5cm}} \underline{\hspace{1.5cm}} \quad TRISTE $=$ \underline{\hspace{1.5cm}} \underline{\hspace{1.5cm}}
  \end{center}

  \item \textbf{Leia e escreva o número de sílabas de cada palavra:}
  \begin{center}
  TREM $\rightarrow$ \underline{\hspace{0.9cm}} \quad TRATOR $\rightarrow$ \underline{\hspace{0.9cm}} \quad TRISTE $\rightarrow$ \underline{\hspace{0.9cm}} \quad TROCA $\rightarrow$ \underline{\hspace{0.9cm}}
  \end{center}

  \item \textbf{Copie as frases em letra cursiva:}
  \begin{center}
  O trem chega cedo.\\
  \underline{\hspace{12cm}}\\
  O trator arou a terra.\\
  \underline{\hspace{12cm}}
  \end{center}

  \item \textbf{Desenhe a frase:}
  \begin{center}
  \begin{tikzpicture}
    \draw[thick, dashed] (0,0) rectangle (12,4);
  \end{tikzpicture}
  \end{center}

\end{enumerate}

\end{exercicio}

%% COMPLEMENTAR C-13: ENCONTRO VR

\plancheta{Folha Complementar C-13 --- Reforço da Folha E-13 (Encontro VR)}

\begin{exercicio}[Folha Complementar C-13 --- Reforço da Folha E-13 (Encontro VR)]

\metadata{C-13}{Silábico-Alfabético}{EF02LP03, EF02LP04, EF02LP07}{D1}{EF02LP04}{CNE/CEB nº 2/2012}

\kumontiming{2}{90}

\noindent\textbf{Nome:} \underline{\hspace{3.2cm}} \quad \textbf{Data:} \underline{\hspace{1.6cm}} \quad \textbf{Nota:} \underline{\hspace{0.9cm}}/10


\begin{enumerate}

  \item \textbf{Copie as palavras em letra cursiva:}
  \begin{center}
  livro \quad palavra \quad lavrar \quad livre
  \end{center}
  \begin{center}
  \underline{\hspace{2.4cm}} \quad \underline{\hspace{2.4cm}} \quad \underline{\hspace{2.4cm}} \quad \underline{\hspace{2.4cm}}
  \end{center}

  \item \textbf{Escreva as sílabas das palavras:}
  \begin{center}
  LIVRO $=$ \underline{\hspace{1.5cm}} \underline{\hspace{1.5cm}} \quad PALAVRA $=$ \underline{\hspace{1.5cm}} \underline{\hspace{1.5cm}} \quad LAVRAR $=$ \underline{\hspace{1.5cm}} \underline{\hspace{1.5cm}}
  \end{center}

  \item \textbf{Leia e escreva o número de sílabas de cada palavra:}
  \begin{center}
  LIVRO $\rightarrow$ \underline{\hspace{0.9cm}} \quad PALAVRA $\rightarrow$ \underline{\hspace{0.9cm}} \quad LAVRAR $\rightarrow$ \underline{\hspace{0.9cm}} \quad LIVRE $\rightarrow$ \underline{\hspace{0.9cm}}
  \end{center}

  \item \textbf{Copie as frases em letra cursiva:}
  \begin{center}
  O livro é meu.\\
  \underline{\hspace{12cm}}\\
  A palavra é longa.\\
  \underline{\hspace{12cm}}
  \end{center}

  \item \textbf{Desenhe a frase:}
  \begin{center}
  \begin{tikzpicture}
    \draw[thick, dashed] (0,0) rectangle (12,4);
  \end{tikzpicture}
  \end{center}

\end{enumerate}

\end{exercicio}

%% COMPLEMENTAR C-14: ENCONTRO BL

\plancheta{Folha Complementar C-14 --- Reforço da Folha E-14 (Encontro BL)}

\begin{exercicio}[Folha Complementar C-14 --- Reforço da Folha E-14 (Encontro BL)]

\metadata{C-14}{Silábico-Alfabético}{EF02LP03, EF02LP04, EF02LP07}{D1}{EF02LP04}{CNE/CEB nº 2/2012}

\kumontiming{2}{90}

\noindent\textbf{Nome:} \underline{\hspace{3.2cm}} \quad \textbf{Data:} \underline{\hspace{1.6cm}} \quad \textbf{Nota:} \underline{\hspace{0.9cm}}/10


\begin{enumerate}

  \item \textbf{Copie as palavras em letra cursiva:}
  \begin{center}
  blusa \quad bloco \quad branco \quad branco
  \end{center}
  \begin{center}
  \underline{\hspace{2.4cm}} \quad \underline{\hspace{2.4cm}} \quad \underline{\hspace{2.4cm}} \quad \underline{\hspace{2.4cm}}
  \end{center}

  \item \textbf{Escreva as sílabas das palavras:}
  \begin{center}
  BLUSA $=$ \underline{\hspace{1.5cm}} \underline{\hspace{1.5cm}} \quad BLOCO $=$ \underline{\hspace{1.5cm}} \underline{\hspace{1.5cm}} \quad BRANCO $=$ \underline{\hspace{1.5cm}} \underline{\hspace{1.5cm}}
  \end{center}

  \item \textbf{Leia e escreva o número de sílabas de cada palavra:}
  \begin{center}
  BLUSA $\rightarrow$ \underline{\hspace{0.9cm}} \quad BLOCO $\rightarrow$ \underline{\hspace{0.9cm}} \quad BRANCO $\rightarrow$ \underline{\hspace{0.9cm}} \quad BRANCO $\rightarrow$ \underline{\hspace{0.9cm}}
  \end{center}

  \item \textbf{Copie as frases em letra cursiva:}
  \begin{center}
  A blusa é azul.\\
  \underline{\hspace{12cm}}\\
  O bloco caiu no chão.\\
  \underline{\hspace{12cm}}
  \end{center}

  \item \textbf{Desenhe a frase:}
  \begin{center}
  \begin{tikzpicture}
    \draw[thick, dashed] (0,0) rectangle (12,4);
  \end{tikzpicture}
  \end{center}

\end{enumerate}

\end{exercicio}

%% COMPLEMENTAR C-15: ENCONTRO CL

\plancheta{Folha Complementar C-15 --- Reforço da Folha E-15 (Encontro CL)}

\begin{exercicio}[Folha Complementar C-15 --- Reforço da Folha E-15 (Encontro CL)]

\metadata{C-15}{Silábico-Alfabético}{EF02LP03, EF02LP04, EF02LP07}{D1}{EF02LP04}{CNE/CEB nº 2/2012}

\kumontiming{2}{90}

\noindent\textbf{Nome:} \underline{\hspace{3.2cm}} \quad \textbf{Data:} \underline{\hspace{1.6cm}} \quad \textbf{Nota:} \underline{\hspace{0.9cm}}/10


\begin{enumerate}

  \item \textbf{Copie as palavras em letra cursiva:}
  \begin{center}
  claro \quad clube \quad clima \quad classe
  \end{center}
  \begin{center}
  \underline{\hspace{2.4cm}} \quad \underline{\hspace{2.4cm}} \quad \underline{\hspace{2.4cm}} \quad \underline{\hspace{2.4cm}}
  \end{center}

  \item \textbf{Escreva as sílabas das palavras:}
  \begin{center}
  CLARO $=$ \underline{\hspace{1.5cm}} \underline{\hspace{1.5cm}} \quad CLUBE $=$ \underline{\hspace{1.5cm}} \underline{\hspace{1.5cm}} \quad CLIMA $=$ \underline{\hspace{1.5cm}} \underline{\hspace{1.5cm}}
  \end{center}

  \item \textbf{Leia e escreva o número de sílabas de cada palavra:}
  \begin{center}
  CLARO $\rightarrow$ \underline{\hspace{0.9cm}} \quad CLUBE $\rightarrow$ \underline{\hspace{0.9cm}} \quad CLIMA $\rightarrow$ \underline{\hspace{0.9cm}} \quad CLASSE $\rightarrow$ \underline{\hspace{0.9cm}}
  \end{center}

  \item \textbf{Copie as frases em letra cursiva:}
  \begin{center}
  O céu está claro.\\
  \underline{\hspace{12cm}}\\
  O clube fica perto.\\
  \underline{\hspace{12cm}}
  \end{center}

  \item \textbf{Desenhe a frase:}
  \begin{center}
  \begin{tikzpicture}
    \draw[thick, dashed] (0,0) rectangle (12,4);
  \end{tikzpicture}
  \end{center}

\end{enumerate}

\end{exercicio}

%% COMPLEMENTAR C-16: ENCONTRO FL

\plancheta{Folha Complementar C-16 --- Reforço da Folha E-16 (Encontro FL)}

\begin{exercicio}[Folha Complementar C-16 --- Reforço da Folha E-16 (Encontro FL)]

\metadata{C-16}{Silábico-Alfabético}{EF02LP03, EF02LP04, EF02LP07}{D1}{EF02LP04}{CNE/CEB nº 2/2012}

\kumontiming{2}{90}

\noindent\textbf{Nome:} \underline{\hspace{3.2cm}} \quad \textbf{Data:} \underline{\hspace{1.6cm}} \quad \textbf{Nota:} \underline{\hspace{0.9cm}}/10


\begin{enumerate}

  \item \textbf{Copie as palavras em letra cursiva:}
  \begin{center}
  flor \quad flauta \quad floco \quad floresta
  \end{center}
  \begin{center}
  \underline{\hspace{2.4cm}} \quad \underline{\hspace{2.4cm}} \quad \underline{\hspace{2.4cm}} \quad \underline{\hspace{2.4cm}}
  \end{center}

  \item \textbf{Escreva as sílabas das palavras:}
  \begin{center}
  FLOR $=$ \underline{\hspace{1.5cm}} \underline{\hspace{1.5cm}} \quad FLAUTA $=$ \underline{\hspace{1.5cm}} \underline{\hspace{1.5cm}} \quad FLOCO $=$ \underline{\hspace{1.5cm}} \underline{\hspace{1.5cm}}
  \end{center}

  \item \textbf{Leia e escreva o número de sílabas de cada palavra:}
  \begin{center}
  FLOR $\rightarrow$ \underline{\hspace{0.9cm}} \quad FLAUTA $\rightarrow$ \underline{\hspace{0.9cm}} \quad FLOCO $\rightarrow$ \underline{\hspace{0.9cm}} \quad FLORESTA $\rightarrow$ \underline{\hspace{0.9cm}}
  \end{center}

  \item \textbf{Copie as frases em letra cursiva:}
  \begin{center}
  A flor abriu no jardim.\\
  \underline{\hspace{12cm}}\\
  A flauta toca doce.\\
  \underline{\hspace{12cm}}
  \end{center}

  \item \textbf{Desenhe a frase:}
  \begin{center}
  \begin{tikzpicture}
    \draw[thick, dashed] (0,0) rectangle (12,4);
  \end{tikzpicture}
  \end{center}

\end{enumerate}

\end{exercicio}

%% COMPLEMENTAR C-17: ENCONTRO GL

\plancheta{Folha Complementar C-17 --- Reforço da Folha E-17 (Encontro GL)}

\begin{exercicio}[Folha Complementar C-17 --- Reforço da Folha E-17 (Encontro GL)]

\metadata{C-17}{Silábico-Alfabético}{EF02LP03, EF02LP04, EF02LP07}{D1}{EF02LP04}{CNE/CEB nº 2/2012}

\kumontiming{2}{90}

\noindent\textbf{Nome:} \underline{\hspace{3.2cm}} \quad \textbf{Data:} \underline{\hspace{1.6cm}} \quad \textbf{Nota:} \underline{\hspace{0.9cm}}/10


\begin{enumerate}

  \item \textbf{Copie as palavras em letra cursiva:}
  \begin{center}
  globo \quad glória \quad glicose \quad inglês
  \end{center}
  \begin{center}
  \underline{\hspace{2.4cm}} \quad \underline{\hspace{2.4cm}} \quad \underline{\hspace{2.4cm}} \quad \underline{\hspace{2.4cm}}
  \end{center}

  \item \textbf{Escreva as sílabas das palavras:}
  \begin{center}
  GLOBO $=$ \underline{\hspace{1.5cm}} \underline{\hspace{1.5cm}} \quad GLÓRIA $=$ \underline{\hspace{1.5cm}} \underline{\hspace{1.5cm}} \quad GLICOSE $=$ \underline{\hspace{1.5cm}} \underline{\hspace{1.5cm}}
  \end{center}

  \item \textbf{Leia e escreva o número de sílabas de cada palavra:}
  \begin{center}
  GLOBO $\rightarrow$ \underline{\hspace{0.9cm}} \quad GLÓRIA $\rightarrow$ \underline{\hspace{0.9cm}} \quad GLICOSE $\rightarrow$ \underline{\hspace{0.9cm}} \quad INGLÊS $\rightarrow$ \underline{\hspace{0.9cm}}
  \end{center}

  \item \textbf{Copie as frases em letra cursiva:}
  \begin{center}
  O globo mostra o mapa.\\
  \underline{\hspace{12cm}}\\
  A glória do time é grande.\\
  \underline{\hspace{12cm}}
  \end{center}

  \item \textbf{Desenhe a frase:}
  \begin{center}
  \begin{tikzpicture}
    \draw[thick, dashed] (0,0) rectangle (12,4);
  \end{tikzpicture}
  \end{center}

\end{enumerate}

\end{exercicio}

%% COMPLEMENTAR C-18: ENCONTRO PL

\plancheta{Folha Complementar C-18 --- Reforço da Folha E-18 (Encontro PL)}

\begin{exercicio}[Folha Complementar C-18 --- Reforço da Folha E-18 (Encontro PL)]

\metadata{C-18}{Silábico-Alfabético}{EF02LP03, EF02LP04, EF02LP07}{D1}{EF02LP04}{CNE/CEB nº 2/2012}

\kumontiming{2}{90}

\noindent\textbf{Nome:} \underline{\hspace{3.2cm}} \quad \textbf{Data:} \underline{\hspace{1.6cm}} \quad \textbf{Nota:} \underline{\hspace{0.9cm}}/10


\begin{enumerate}

  \item \textbf{Copie as palavras em letra cursiva:}
  \begin{center}
  placa \quad planta \quad plano \quad plástico
  \end{center}
  \begin{center}
  \underline{\hspace{2.4cm}} \quad \underline{\hspace{2.4cm}} \quad \underline{\hspace{2.4cm}} \quad \underline{\hspace{2.4cm}}
  \end{center}

  \item \textbf{Escreva as sílabas das palavras:}
  \begin{center}
  PLACA $=$ \underline{\hspace{1.5cm}} \underline{\hspace{1.5cm}} \quad PLANTA $=$ \underline{\hspace{1.5cm}} \underline{\hspace{1.5cm}} \quad PLANO $=$ \underline{\hspace{1.5cm}} \underline{\hspace{1.5cm}}
  \end{center}

  \item \textbf{Leia e escreva o número de sílabas de cada palavra:}
  \begin{center}
  PLACA $\rightarrow$ \underline{\hspace{0.9cm}} \quad PLANTA $\rightarrow$ \underline{\hspace{0.9cm}} \quad PLANO $\rightarrow$ \underline{\hspace{0.9cm}} \quad PLÁSTICO $\rightarrow$ \underline{\hspace{0.9cm}}
  \end{center}

  \item \textbf{Copie as frases em letra cursiva:}
  \begin{center}
  A placa indica o caminho.\\
  \underline{\hspace{12cm}}\\
  A planta cresceu.\\
  \underline{\hspace{12cm}}
  \end{center}

  \item \textbf{Desenhe a frase:}
  \begin{center}
  \begin{tikzpicture}
    \draw[thick, dashed] (0,0) rectangle (12,4);
  \end{tikzpicture}
  \end{center}

\end{enumerate}

\end{exercicio}

%% COMPLEMENTAR C-19: DÍGRAFO RR

\plancheta{Folha Complementar C-19 --- Reforço da Folha E-19 (Dígrafo RR)}

\begin{exercicio}[Folha Complementar C-19 --- Reforço da Folha E-19 (Dígrafo RR)]

\metadata{C-19}{Silábico-Alfabético}{EF02LP03, EF02LP04, EF02LP07}{D1}{EF02LP04}{CNE/CEB nº 2/2012}

\kumontiming{2}{90}

\noindent\textbf{Nome:} \underline{\hspace{3.2cm}} \quad \textbf{Data:} \underline{\hspace{1.6cm}} \quad \textbf{Nota:} \underline{\hspace{0.9cm}}/10


\begin{enumerate}

  \item \textbf{Copie as palavras em letra cursiva:}
  \begin{center}
  carro \quad terra \quad barro \quad sorriso
  \end{center}
  \begin{center}
  \underline{\hspace{2.4cm}} \quad \underline{\hspace{2.4cm}} \quad \underline{\hspace{2.4cm}} \quad \underline{\hspace{2.4cm}}
  \end{center}

  \item \textbf{Escreva as sílabas das palavras:}
  \begin{center}
  CARRO $=$ \underline{\hspace{1.5cm}} \underline{\hspace{1.5cm}} \quad TERRA $=$ \underline{\hspace{1.5cm}} \underline{\hspace{1.5cm}} \quad BARRO $=$ \underline{\hspace{1.5cm}} \underline{\hspace{1.5cm}}
  \end{center}

  \item \textbf{Leia e escreva o número de sílabas de cada palavra:}
  \begin{center}
  CARRO $\rightarrow$ \underline{\hspace{0.9cm}} \quad TERRA $\rightarrow$ \underline{\hspace{0.9cm}} \quad BARRO $\rightarrow$ \underline{\hspace{0.9cm}} \quad SORRISO $\rightarrow$ \underline{\hspace{0.9cm}}
  \end{center}

  \item \textbf{Copie as frases em letra cursiva:}
  \begin{center}
  O carro andou na terra.\\
  \underline{\hspace{12cm}}\\
  O barro sujou o sapato.\\
  \underline{\hspace{12cm}}
  \end{center}

  \item \textbf{Desenhe a frase:}
  \begin{center}
  \begin{tikzpicture}
    \draw[thick, dashed] (0,0) rectangle (12,4);
  \end{tikzpicture}
  \end{center}

\end{enumerate}

\end{exercicio}

%% COMPLEMENTAR C-20: DÍGRAFO SS

\plancheta{Folha Complementar C-20 --- Reforço da Folha E-20 (Dígrafo SS)}

\begin{exercicio}[Folha Complementar C-20 --- Reforço da Folha E-20 (Dígrafo SS)]

\metadata{C-20}{Silábico-Alfabético}{EF02LP03, EF02LP04, EF02LP07}{D1}{EF02LP04}{CNE/CEB nº 2/2012}

\kumontiming{2}{90}

\noindent\textbf{Nome:} \underline{\hspace{3.2cm}} \quad \textbf{Data:} \underline{\hspace{1.6cm}} \quad \textbf{Nota:} \underline{\hspace{0.9cm}}/10


\begin{enumerate}

  \item \textbf{Copie as palavras em letra cursiva:}
  \begin{center}
  passar \quad osso \quad assado \quad pássaro
  \end{center}
  \begin{center}
  \underline{\hspace{2.4cm}} \quad \underline{\hspace{2.4cm}} \quad \underline{\hspace{2.4cm}} \quad \underline{\hspace{2.4cm}}
  \end{center}

  \item \textbf{Escreva as sílabas das palavras:}
  \begin{center}
  PASSAR $=$ \underline{\hspace{1.5cm}} \underline{\hspace{1.5cm}} \quad OSSO $=$ \underline{\hspace{1.5cm}} \underline{\hspace{1.5cm}} \quad ASSADO $=$ \underline{\hspace{1.5cm}} \underline{\hspace{1.5cm}}
  \end{center}

  \item \textbf{Leia e escreva o número de sílabas de cada palavra:}
  \begin{center}
  PASSAR $\rightarrow$ \underline{\hspace{0.9cm}} \quad OSSO $\rightarrow$ \underline{\hspace{0.9cm}} \quad ASSADO $\rightarrow$ \underline{\hspace{0.9cm}} \quad PÁSSARO $\rightarrow$ \underline{\hspace{0.9cm}}
  \end{center}

  \item \textbf{Copie as frases em letra cursiva:}
  \begin{center}
  O pássaro cantou cedo.\\
  \underline{\hspace{12cm}}\\
  O osso é do cachorro.\\
  \underline{\hspace{12cm}}
  \end{center}

  \item \textbf{Desenhe a frase:}
  \begin{center}
  \begin{tikzpicture}
    \draw[thick, dashed] (0,0) rectangle (12,4);
  \end{tikzpicture}
  \end{center}

\end{enumerate}

\end{exercicio}

%% COMPLEMENTAR C-21: S ENTRE VOGAIS E Z

\plancheta{Folha Complementar C-21 --- Reforço da Folha E-21 (S entre vogais e Z)}

\begin{exercicio}[Folha Complementar C-21 --- Reforço da Folha E-21 (S entre vogais e Z)]

\metadata{C-21}{Silábico-Alfabético}{EF02LP03, EF02LP04, EF02LP07}{D1}{EF02LP04}{CNE/CEB nº 2/2012}

\kumontiming{2}{90}

\noindent\textbf{Nome:} \underline{\hspace{3.2cm}} \quad \textbf{Data:} \underline{\hspace{1.6cm}} \quad \textbf{Nota:} \underline{\hspace{0.9cm}}/10


\begin{enumerate}

  \item \textbf{Copie as palavras em letra cursiva:}
  \begin{center}
  casa \quad mesa \quad asa \quad poesia
  \end{center}
  \begin{center}
  \underline{\hspace{2.4cm}} \quad \underline{\hspace{2.4cm}} \quad \underline{\hspace{2.4cm}} \quad \underline{\hspace{2.4cm}}
  \end{center}

  \item \textbf{Escreva as sílabas das palavras:}
  \begin{center}
  CASA $=$ \underline{\hspace{1.5cm}} \underline{\hspace{1.5cm}} \quad MESA $=$ \underline{\hspace{1.5cm}} \underline{\hspace{1.5cm}} \quad ASA $=$ \underline{\hspace{1.5cm}} \underline{\hspace{1.5cm}}
  \end{center}

  \item \textbf{Leia e escreva o número de sílabas de cada palavra:}
  \begin{center}
  CASA $\rightarrow$ \underline{\hspace{0.9cm}} \quad MESA $\rightarrow$ \underline{\hspace{0.9cm}} \quad ASA $\rightarrow$ \underline{\hspace{0.9cm}} \quad POESIA $\rightarrow$ \underline{\hspace{0.9cm}}
  \end{center}

  \item \textbf{Copie as frases em letra cursiva:}
  \begin{center}
  A casa é na rua da escola.\\
  \underline{\hspace{12cm}}\\
  A mesa tem quatro cadeiras.\\
  \underline{\hspace{12cm}}
  \end{center}

  \item \textbf{Desenhe a frase:}
  \begin{center}
  \begin{tikzpicture}
    \draw[thick, dashed] (0,0) rectangle (12,4);
  \end{tikzpicture}
  \end{center}

\end{enumerate}

\end{exercicio}

%% COMPLEMENTAR C-22: LETRA Ç

\plancheta{Folha Complementar C-22 --- Reforço da Folha E-22 (Letra Ç)}

\begin{exercicio}[Folha Complementar C-22 --- Reforço da Folha E-22 (Letra Ç)]

\metadata{C-22}{Silábico-Alfabético}{EF02LP03, EF02LP04, EF02LP07}{D1}{EF02LP04}{CNE/CEB nº 2/2012}

\kumontiming{2}{90}

\noindent\textbf{Nome:} \underline{\hspace{3.2cm}} \quad \textbf{Data:} \underline{\hspace{1.6cm}} \quad \textbf{Nota:} \underline{\hspace{0.9cm}}/10


\begin{enumerate}

  \item \textbf{Copie as palavras em letra cursiva:}
  \begin{center}
  coração \quad caçarola \quad açúcar \quad poço
  \end{center}
  \begin{center}
  \underline{\hspace{2.4cm}} \quad \underline{\hspace{2.4cm}} \quad \underline{\hspace{2.4cm}} \quad \underline{\hspace{2.4cm}}
  \end{center}

  \item \textbf{Escreva as sílabas das palavras:}
  \begin{center}
  CORAÇÃO $=$ \underline{\hspace{1.5cm}} \underline{\hspace{1.5cm}} \quad CAÇAROLA $=$ \underline{\hspace{1.5cm}} \underline{\hspace{1.5cm}} \quad AÇÚCAR $=$ \underline{\hspace{1.5cm}} \underline{\hspace{1.5cm}}
  \end{center}

  \item \textbf{Leia e escreva o número de sílabas de cada palavra:}
  \begin{center}
  CORAÇÃO $\rightarrow$ \underline{\hspace{0.9cm}} \quad CAÇAROLA $\rightarrow$ \underline{\hspace{0.9cm}} \quad AÇÚCAR $\rightarrow$ \underline{\hspace{0.9cm}} \quad POÇO $\rightarrow$ \underline{\hspace{0.9cm}}
  \end{center}

  \item \textbf{Copie as frases em letra cursiva:}
  \begin{center}
  O coração bate forte.\\
  \underline{\hspace{12cm}}\\
  O açúcar adoça o café.\\
  \underline{\hspace{12cm}}
  \end{center}

  \item \textbf{Desenhe a frase:}
  \begin{center}
  \begin{tikzpicture}
    \draw[thick, dashed] (0,0) rectangle (12,4);
  \end{tikzpicture}
  \end{center}

\end{enumerate}

\end{exercicio}

%% COMPLEMENTAR C-23: SÍLABAS AM

\plancheta{Folha Complementar C-23 --- Reforço da Folha E-23 (Sílabas AM)}

\begin{exercicio}[Folha Complementar C-23 --- Reforço da Folha E-23 (Sílabas AM)]

\metadata{C-23}{Silábico-Alfabético}{EF02LP03, EF02LP04, EF02LP07}{D1}{EF02LP04}{CNE/CEB nº 2/2012}

\kumontiming{2}{90}

\noindent\textbf{Nome:} \underline{\hspace{3.2cm}} \quad \textbf{Data:} \underline{\hspace{1.6cm}} \quad \textbf{Nota:} \underline{\hspace{0.9cm}}/10


\begin{enumerate}

  \item \textbf{Copie as palavras em letra cursiva:}
  \begin{center}
  também \quad sambar \quad cama \quad chamar
  \end{center}
  \begin{center}
  \underline{\hspace{2.4cm}} \quad \underline{\hspace{2.4cm}} \quad \underline{\hspace{2.4cm}} \quad \underline{\hspace{2.4cm}}
  \end{center}

  \item \textbf{Escreva as sílabas das palavras:}
  \begin{center}
  TAMBÉM $=$ \underline{\hspace{1.5cm}} \underline{\hspace{1.5cm}} \quad SAMBAR $=$ \underline{\hspace{1.5cm}} \underline{\hspace{1.5cm}} \quad CAMA $=$ \underline{\hspace{1.5cm}} \underline{\hspace{1.5cm}}
  \end{center}

  \item \textbf{Complete a tabela de sílabas:}

\begin{center}
\resizebox{\textwidth}{!}{%
\begin{tabular}{ccccc}
  \sil{AM}{A} & \sil{AM}{E} & \sil{AM}{I} & \sil{AM}{O} & \sil{AM}{U} \\
  \end{tabular}}
\end{center}

  \item \textbf{Copie as frases em letra cursiva:}
  \begin{center}
  Eu também quero lanche.\\
  \underline{\hspace{12cm}}\\
  As crianças brincam no pátio.\\
  \underline{\hspace{12cm}}
  \end{center}

  \item \textbf{Desenhe a frase:}
  \begin{center}
  \begin{tikzpicture}
    \draw[thick, dashed] (0,0) rectangle (12,4);
  \end{tikzpicture}
  \end{center}

\end{enumerate}

\end{exercicio}

%% COMPLEMENTAR C-24: SÍLABAS AN

\plancheta{Folha Complementar C-24 --- Reforço da Folha E-24 (Sílabas AN)}

\begin{exercicio}[Folha Complementar C-24 --- Reforço da Folha E-24 (Sílabas AN)]

\metadata{C-24}{Silábico-Alfabético}{EF02LP03, EF02LP04, EF02LP07}{D1}{EF02LP04}{CNE/CEB nº 2/2012}

\kumontiming{2}{90}

\noindent\textbf{Nome:} \underline{\hspace{3.2cm}} \quad \textbf{Data:} \underline{\hspace{1.6cm}} \quad \textbf{Nota:} \underline{\hspace{0.9cm}}/10


\begin{enumerate}

  \item \textbf{Copie as palavras em letra cursiva:}
  \begin{center}
  canto \quad sangue \quad banco \quad campo
  \end{center}
  \begin{center}
  \underline{\hspace{2.4cm}} \quad \underline{\hspace{2.4cm}} \quad \underline{\hspace{2.4cm}} \quad \underline{\hspace{2.4cm}}
  \end{center}

  \item \textbf{Escreva as sílabas das palavras:}
  \begin{center}
  CANTO $=$ \underline{\hspace{1.5cm}} \underline{\hspace{1.5cm}} \quad SANGUE $=$ \underline{\hspace{1.5cm}} \underline{\hspace{1.5cm}} \quad BANCO $=$ \underline{\hspace{1.5cm}} \underline{\hspace{1.5cm}}
  \end{center}

  \item \textbf{Complete a tabela de sílabas:}

\begin{center}
\resizebox{\textwidth}{!}{%
\begin{tabular}{ccccc}
  \sil{AN}{A} & \sil{AN}{E} & \sil{AN}{I} & \sil{AN}{O} & \sil{AN}{U} \\
  \end{tabular}}
\end{center}

  \item \textbf{Copie as frases em letra cursiva:}
  \begin{center}
  O canto da sala é azul.\\
  \underline{\hspace{12cm}}\\
  O banco está na praça.\\
  \underline{\hspace{12cm}}
  \end{center}

  \item \textbf{Desenhe a frase:}
  \begin{center}
  \begin{tikzpicture}
    \draw[thick, dashed] (0,0) rectangle (12,4);
  \end{tikzpicture}
  \end{center}

\end{enumerate}

\end{exercicio}

%% COMPLEMENTAR C-25: X COM SOM DE CH

\plancheta{Folha Complementar C-25 --- Reforço da Folha E-25 (X com som de CH)}

\begin{exercicio}[Folha Complementar C-25 --- Reforço da Folha E-25 (X com som de CH)]

\metadata{C-25}{Silábico-Alfabético}{EF02LP03, EF02LP04, EF02LP07}{D1}{EF02LP04}{CNE/CEB nº 2/2012}

\kumontiming{2}{90}

\noindent\textbf{Nome:} \underline{\hspace{3.2cm}} \quad \textbf{Data:} \underline{\hspace{1.6cm}} \quad \textbf{Nota:} \underline{\hspace{0.9cm}}/10


\begin{enumerate}

  \item \textbf{Copie as palavras em letra cursiva:}
  \begin{center}
  xadrez \quad peixe \quad caixa \quad xícara
  \end{center}
  \begin{center}
  \underline{\hspace{2.4cm}} \quad \underline{\hspace{2.4cm}} \quad \underline{\hspace{2.4cm}} \quad \underline{\hspace{2.4cm}}
  \end{center}

  \item \textbf{Escreva as sílabas das palavras:}
  \begin{center}
  XADREZ $=$ \underline{\hspace{1.5cm}} \underline{\hspace{1.5cm}} \quad PEIXE $=$ \underline{\hspace{1.5cm}} \underline{\hspace{1.5cm}} \quad CAIXA $=$ \underline{\hspace{1.5cm}} \underline{\hspace{1.5cm}}
  \end{center}

  \item \textbf{Leia e escreva o número de sílabas de cada palavra:}
  \begin{center}
  XADREZ $\rightarrow$ \underline{\hspace{0.9cm}} \quad PEIXE $\rightarrow$ \underline{\hspace{0.9cm}} \quad CAIXA $\rightarrow$ \underline{\hspace{0.9cm}} \quad XÍCARA $\rightarrow$ \underline{\hspace{0.9cm}}
  \end{center}

  \item \textbf{Copie as frases em letra cursiva:}
  \begin{center}
  O peixe nada no aquário.\\
  \underline{\hspace{12cm}}\\
  A caixa tem brinquedos.\\
  \underline{\hspace{12cm}}
  \end{center}

  \item \textbf{Desenhe a frase:}
  \begin{center}
  \begin{tikzpicture}
    \draw[thick, dashed] (0,0) rectangle (12,4);
  \end{tikzpicture}
  \end{center}

\end{enumerate}

\end{exercicio}

%% COMPLEMENTAR C-26: X COM SOM DE Z

\plancheta{Folha Complementar C-26 --- Reforço da Folha E-26 (X com som de Z)}

\begin{exercicio}[Folha Complementar C-26 --- Reforço da Folha E-26 (X com som de Z)]

\metadata{C-26}{Silábico-Alfabético}{EF02LP03, EF02LP04, EF02LP07}{D1}{EF02LP04}{CNE/CEB nº 2/2012}

\kumontiming{2}{90}

\noindent\textbf{Nome:} \underline{\hspace{3.2cm}} \quad \textbf{Data:} \underline{\hspace{1.6cm}} \quad \textbf{Nota:} \underline{\hspace{0.9cm}}/10


\begin{enumerate}

  \item \textbf{Copie as palavras em letra cursiva:}
  \begin{center}
  exame \quad exemplo \quad exato \quad executar
  \end{center}
  \begin{center}
  \underline{\hspace{2.4cm}} \quad \underline{\hspace{2.4cm}} \quad \underline{\hspace{2.4cm}} \quad \underline{\hspace{2.4cm}}
  \end{center}

  \item \textbf{Escreva as sílabas das palavras:}
  \begin{center}
  EXAME $=$ \underline{\hspace{1.5cm}} \underline{\hspace{1.5cm}} \quad EXEMPLO $=$ \underline{\hspace{1.5cm}} \underline{\hspace{1.5cm}} \quad EXATO $=$ \underline{\hspace{1.5cm}} \underline{\hspace{1.5cm}}
  \end{center}

  \item \textbf{Leia e escreva o número de sílabas de cada palavra:}
  \begin{center}
  EXAME $\rightarrow$ \underline{\hspace{0.9cm}} \quad EXEMPLO $\rightarrow$ \underline{\hspace{0.9cm}} \quad EXATO $\rightarrow$ \underline{\hspace{0.9cm}} \quad EXECUTAR $\rightarrow$ \underline{\hspace{0.9cm}}
  \end{center}

  \item \textbf{Copie as frases em letra cursiva:}
  \begin{center}
  O exame foi fácil.\\
  \underline{\hspace{12cm}}\\
  O exemplo ajudou a turma.\\
  \underline{\hspace{12cm}}
  \end{center}

  \item \textbf{Desenhe a frase:}
  \begin{center}
  \begin{tikzpicture}
    \draw[thick, dashed] (0,0) rectangle (12,4);
  \end{tikzpicture}
  \end{center}

\end{enumerate}

\end{exercicio}

%% COMPLEMENTAR C-27: X COM SOM DE KS

\plancheta{Folha Complementar C-27 --- Reforço da Folha E-27 (X com som de KS)}

\begin{exercicio}[Folha Complementar C-27 --- Reforço da Folha E-27 (X com som de KS)]

\metadata{C-27}{Silábico-Alfabético}{EF02LP03, EF02LP04, EF02LP07}{D1}{EF02LP04}{CNE/CEB nº 2/2012}

\kumontiming{2}{90}

\noindent\textbf{Nome:} \underline{\hspace{3.2cm}} \quad \textbf{Data:} \underline{\hspace{1.6cm}} \quad \textbf{Nota:} \underline{\hspace{0.9cm}}/10


\begin{enumerate}

  \item \textbf{Copie as palavras em letra cursiva:}
  \begin{center}
  táxi \quad fixo \quad saxofone \quad oxigênio
  \end{center}
  \begin{center}
  \underline{\hspace{2.4cm}} \quad \underline{\hspace{2.4cm}} \quad \underline{\hspace{2.4cm}} \quad \underline{\hspace{2.4cm}}
  \end{center}

  \item \textbf{Escreva as sílabas das palavras:}
  \begin{center}
  TÁXI $=$ \underline{\hspace{1.5cm}} \underline{\hspace{1.5cm}} \quad FIXO $=$ \underline{\hspace{1.5cm}} \underline{\hspace{1.5cm}} \quad SAXOFONE $=$ \underline{\hspace{1.5cm}} \underline{\hspace{1.5cm}}
  \end{center}

  \item \textbf{Leia e escreva o número de sílabas de cada palavra:}
  \begin{center}
  TÁXI $\rightarrow$ \underline{\hspace{0.9cm}} \quad FIXO $\rightarrow$ \underline{\hspace{0.9cm}} \quad SAXOFONE $\rightarrow$ \underline{\hspace{0.9cm}} \quad OXIGÊNIO $\rightarrow$ \underline{\hspace{0.9cm}}
  \end{center}

  \item \textbf{Copie as frases em letra cursiva:}
  \begin{center}
  O táxi parou na esquina.\\
  \underline{\hspace{12cm}}\\
  O saxofone toca jazz.\\
  \underline{\hspace{12cm}}
  \end{center}

  \item \textbf{Desenhe a frase:}
  \begin{center}
  \begin{tikzpicture}
    \draw[thick, dashed] (0,0) rectangle (12,4);
  \end{tikzpicture}
  \end{center}

\end{enumerate}

\end{exercicio}

%% COMPLEMENTAR C-28: REVISÃO DAS GRAFIAS DO ANO

\plancheta{Folha Complementar C-28 --- Reforço da Folha E-28 (Revisão das Grafias do Ano)}

\begin{exercicio}[Folha Complementar C-28 --- Reforço da Folha E-28 (Revisão das Grafias do Ano)]

\metadata{C-28}{Silábico-Alfabético}{EF02LP03, EF02LP04, EF02LP07}{D1}{EF02LP04}{CNE/CEB nº 2/2012}

\kumontiming{2}{90}

\noindent\textbf{Nome:} \underline{\hspace{3.2cm}} \quad \textbf{Data:} \underline{\hspace{1.6cm}} \quad \textbf{Nota:} \underline{\hspace{0.9cm}}/10


\begin{enumerate}

  \item \textbf{Copie as palavras em letra cursiva:}
  \begin{center}
  chave \quad coelho \quad ninho \quad queijo
  \end{center}
  \begin{center}
  \underline{\hspace{2.4cm}} \quad \underline{\hspace{2.4cm}} \quad \underline{\hspace{2.4cm}} \quad \underline{\hspace{2.4cm}}
  \end{center}

  \item \textbf{Escreva as sílabas das palavras:}
  \begin{center}
  CHAVE $=$ \underline{\hspace{1.5cm}} \underline{\hspace{1.5cm}} \quad COELHO $=$ \underline{\hspace{1.5cm}} \underline{\hspace{1.5cm}} \quad NINHO $=$ \underline{\hspace{1.5cm}} \underline{\hspace{1.5cm}}
  \end{center}

  \item \textbf{Leia e escreva o número de sílabas de cada palavra:}
  \begin{center}
  CHAVE $\rightarrow$ \underline{\hspace{0.9cm}} \quad COELHO $\rightarrow$ \underline{\hspace{0.9cm}} \quad NINHO $\rightarrow$ \underline{\hspace{0.9cm}} \quad QUEIJO $\rightarrow$ \underline{\hspace{0.9cm}}
  \end{center}

  \item \textbf{Copie as frases em letra cursiva:}
  \begin{center}
  A chave abriu o cofre.\\
  \underline{\hspace{12cm}}\\
  O coelho correu no jardim.\\
  \underline{\hspace{12cm}}
  \end{center}

  \item \textbf{Desenhe a frase:}
  \begin{center}
  \begin{tikzpicture}
    \draw[thick, dashed] (0,0) rectangle (12,4);
  \end{tikzpicture}
  \end{center}

\end{enumerate}

\end{exercicio}

%% COMPLEMENTAR C-29: REVISÃO CH, LH E NH

\plancheta{Folha Complementar C-29 --- Reforço da Folha E-29 (Revisão CH, LH e NH)}

\begin{exercicio}[Folha Complementar C-29 --- Reforço da Folha E-29 (Revisão CH, LH e NH)]

\metadata{C-29}{Silábico-Alfabético}{EF02LP03, EF02LP04, EF02LP07}{D1}{EF02LP04}{CNE/CEB nº 2/2012}

\kumontiming{2}{90}

\noindent\textbf{Nome:} \underline{\hspace{3.2cm}} \quad \textbf{Data:} \underline{\hspace{1.6cm}} \quad \textbf{Nota:} \underline{\hspace{0.9cm}}/10


\begin{enumerate}

  \item \textbf{Copie as palavras em letra cursiva:}
  \begin{center}
  chuva \quad coelho \quad ninho \quad chave
  \end{center}
  \begin{center}
  \underline{\hspace{2.4cm}} \quad \underline{\hspace{2.4cm}} \quad \underline{\hspace{2.4cm}} \quad \underline{\hspace{2.4cm}}
  \end{center}

  \item \textbf{Escreva as sílabas das palavras:}
  \begin{center}
  CHUVA $=$ \underline{\hspace{1.5cm}} \underline{\hspace{1.5cm}} \quad COELHO $=$ \underline{\hspace{1.5cm}} \underline{\hspace{1.5cm}} \quad NINHO $=$ \underline{\hspace{1.5cm}} \underline{\hspace{1.5cm}}
  \end{center}

  \item \textbf{Leia e escreva o número de sílabas de cada palavra:}
  \begin{center}
  CHUVA $\rightarrow$ \underline{\hspace{0.9cm}} \quad COELHO $\rightarrow$ \underline{\hspace{0.9cm}} \quad NINHO $\rightarrow$ \underline{\hspace{0.9cm}} \quad CHAVE $\rightarrow$ \underline{\hspace{0.9cm}}
  \end{center}

  \item \textbf{Copie as frases em letra cursiva:}
  \begin{center}
  A chuva molhou o ninho.\\
  \underline{\hspace{12cm}}\\
  O coelho escondeu-se na folha.\\
  \underline{\hspace{12cm}}
  \end{center}

  \item \textbf{Desenhe a frase:}
  \begin{center}
  \begin{tikzpicture}
    \draw[thick, dashed] (0,0) rectangle (12,4);
  \end{tikzpicture}
  \end{center}

\end{enumerate}

\end{exercicio}

%% COMPLEMENTAR C-30: REVISÃO QUE, QUI, GUE, GUI

\plancheta{Folha Complementar C-30 --- Reforço da Folha E-30 (Revisão QUE, QUI, GUE, GUI)}

\begin{exercicio}[Folha Complementar C-30 --- Reforço da Folha E-30 (Revisão QUE, QUI, GUE, GUI)]

\metadata{C-30}{Silábico-Alfabético}{EF02LP03, EF02LP04, EF02LP07}{D1}{EF02LP04}{CNE/CEB nº 2/2012}

\kumontiming{2}{90}

\noindent\textbf{Nome:} \underline{\hspace{3.2cm}} \quad \textbf{Data:} \underline{\hspace{1.6cm}} \quad \textbf{Nota:} \underline{\hspace{0.9cm}}/10


\begin{enumerate}

  \item \textbf{Copie as palavras em letra cursiva:}
  \begin{center}
  queijo \quad quilo \quad quente \quad guerra
  \end{center}
  \begin{center}
  \underline{\hspace{2.4cm}} \quad \underline{\hspace{2.4cm}} \quad \underline{\hspace{2.4cm}} \quad \underline{\hspace{2.4cm}}
  \end{center}

  \item \textbf{Escreva as sílabas das palavras:}
  \begin{center}
  QUEIJO $=$ \underline{\hspace{1.5cm}} \underline{\hspace{1.5cm}} \quad QUILO $=$ \underline{\hspace{1.5cm}} \underline{\hspace{1.5cm}} \quad QUENTE $=$ \underline{\hspace{1.5cm}} \underline{\hspace{1.5cm}}
  \end{center}

  \item \textbf{Leia e escreva o número de sílabas de cada palavra:}
  \begin{center}
  QUEIJO $\rightarrow$ \underline{\hspace{0.9cm}} \quad QUILO $\rightarrow$ \underline{\hspace{0.9cm}} \quad QUENTE $\rightarrow$ \underline{\hspace{0.9cm}} \quad GUERRA $\rightarrow$ \underline{\hspace{0.9cm}}
  \end{center}

  \item \textbf{Copie as frases em letra cursiva:}
  \begin{center}
  O queijo está com o periquito.\\
  \underline{\hspace{12cm}}\\
  Eu quero um quilo de água.\\
  \underline{\hspace{12cm}}
  \end{center}

  \item \textbf{Desenhe a frase:}
  \begin{center}
  \begin{tikzpicture}
    \draw[thick, dashed] (0,0) rectangle (12,4);
  \end{tikzpicture}
  \end{center}

\end{enumerate}

\end{exercicio}

%% COMPLEMENTAR C-31: REVISÃO ENCONTROS 1 (B–G)

\plancheta{Folha Complementar C-31 --- Reforço da Folha E-31 (Revisão Encontros 1 (B–G))}

\begin{exercicio}[Folha Complementar C-31 --- Reforço da Folha E-31 (Revisão Encontros 1 (B–G))]

\metadata{C-31}{Silábico-Alfabético}{EF02LP03, EF02LP04, EF02LP07}{D1}{EF02LP04}{CNE/CEB nº 2/2012}

\kumontiming{2}{90}

\noindent\textbf{Nome:} \underline{\hspace{3.2cm}} \quad \textbf{Data:} \underline{\hspace{1.6cm}} \quad \textbf{Nota:} \underline{\hspace{0.9cm}}/10


\begin{enumerate}

  \item \textbf{Copie as palavras em letra cursiva:}
  \begin{center}
  braço \quad cravo \quad dragão \quad fruta
  \end{center}
  \begin{center}
  \underline{\hspace{2.4cm}} \quad \underline{\hspace{2.4cm}} \quad \underline{\hspace{2.4cm}} \quad \underline{\hspace{2.4cm}}
  \end{center}

  \item \textbf{Escreva as sílabas das palavras:}
  \begin{center}
  BRAÇO $=$ \underline{\hspace{1.5cm}} \underline{\hspace{1.5cm}} \quad CRAVO $=$ \underline{\hspace{1.5cm}} \underline{\hspace{1.5cm}} \quad DRAGÃO $=$ \underline{\hspace{1.5cm}} \underline{\hspace{1.5cm}}
  \end{center}

  \item \textbf{Leia e escreva o número de sílabas de cada palavra:}
  \begin{center}
  BRAÇO $\rightarrow$ \underline{\hspace{0.9cm}} \quad CRAVO $\rightarrow$ \underline{\hspace{0.9cm}} \quad DRAGÃO $\rightarrow$ \underline{\hspace{0.9cm}} \quad FRUTA $\rightarrow$ \underline{\hspace{0.9cm}}
  \end{center}

  \item \textbf{Copie as frases em letra cursiva:}
  \begin{center}
  O braço segura o cravo.\\
  \underline{\hspace{12cm}}\\
  O dragão comeu a fruta.\\
  \underline{\hspace{12cm}}
  \end{center}

  \item \textbf{Desenhe a frase:}
  \begin{center}
  \begin{tikzpicture}
    \draw[thick, dashed] (0,0) rectangle (12,4);
  \end{tikzpicture}
  \end{center}

\end{enumerate}

\end{exercicio}

%% COMPLEMENTAR C-32: REVISÃO ENCONTROS 2 (L)

\plancheta{Folha Complementar C-32 --- Reforço da Folha E-32 (Revisão Encontros 2 (L))}

\begin{exercicio}[Folha Complementar C-32 --- Reforço da Folha E-32 (Revisão Encontros 2 (L))]

\metadata{C-32}{Silábico-Alfabético}{EF02LP03, EF02LP04, EF02LP07}{D1}{EF02LP04}{CNE/CEB nº 2/2012}

\kumontiming{2}{90}

\noindent\textbf{Nome:} \underline{\hspace{3.2cm}} \quad \textbf{Data:} \underline{\hspace{1.6cm}} \quad \textbf{Nota:} \underline{\hspace{0.9cm}}/10


\begin{enumerate}

  \item \textbf{Copie as palavras em letra cursiva:}
  \begin{center}
  blusa \quad claro \quad flor \quad globo
  \end{center}
  \begin{center}
  \underline{\hspace{2.4cm}} \quad \underline{\hspace{2.4cm}} \quad \underline{\hspace{2.4cm}} \quad \underline{\hspace{2.4cm}}
  \end{center}

  \item \textbf{Escreva as sílabas das palavras:}
  \begin{center}
  BLUSA $=$ \underline{\hspace{1.5cm}} \underline{\hspace{1.5cm}} \quad CLARO $=$ \underline{\hspace{1.5cm}} \underline{\hspace{1.5cm}} \quad FLOR $=$ \underline{\hspace{1.5cm}} \underline{\hspace{1.5cm}}
  \end{center}

  \item \textbf{Leia e escreva o número de sílabas de cada palavra:}
  \begin{center}
  BLUSA $\rightarrow$ \underline{\hspace{0.9cm}} \quad CLARO $\rightarrow$ \underline{\hspace{0.9cm}} \quad FLOR $\rightarrow$ \underline{\hspace{0.9cm}} \quad GLOBO $\rightarrow$ \underline{\hspace{0.9cm}}
  \end{center}

  \item \textbf{Copie as frases em letra cursiva:}
  \begin{center}
  A blusa clara tem flor.\\
  \underline{\hspace{12cm}}\\
  O globo está sobre o atlas.\\
  \underline{\hspace{12cm}}
  \end{center}

  \item \textbf{Desenhe a frase:}
  \begin{center}
  \begin{tikzpicture}
    \draw[thick, dashed] (0,0) rectangle (12,4);
  \end{tikzpicture}
  \end{center}

\end{enumerate}

\end{exercicio}

%% COMPLEMENTAR C-33: REVISÃO RR E SS

\plancheta{Folha Complementar C-33 --- Reforço da Folha E-33 (Revisão RR e SS)}

\begin{exercicio}[Folha Complementar C-33 --- Reforço da Folha E-33 (Revisão RR e SS)]

\metadata{C-33}{Silábico-Alfabético}{EF02LP03, EF02LP04, EF02LP07}{D1}{EF02LP04}{CNE/CEB nº 2/2012}

\kumontiming{2}{90}

\noindent\textbf{Nome:} \underline{\hspace{3.2cm}} \quad \textbf{Data:} \underline{\hspace{1.6cm}} \quad \textbf{Nota:} \underline{\hspace{0.9cm}}/10


\begin{enumerate}

  \item \textbf{Copie as palavras em letra cursiva:}
  \begin{center}
  carro \quad terra \quad barro \quad sorriso
  \end{center}
  \begin{center}
  \underline{\hspace{2.4cm}} \quad \underline{\hspace{2.4cm}} \quad \underline{\hspace{2.4cm}} \quad \underline{\hspace{2.4cm}}
  \end{center}

  \item \textbf{Escreva as sílabas das palavras:}
  \begin{center}
  CARRO $=$ \underline{\hspace{1.5cm}} \underline{\hspace{1.5cm}} \quad TERRA $=$ \underline{\hspace{1.5cm}} \underline{\hspace{1.5cm}} \quad BARRO $=$ \underline{\hspace{1.5cm}} \underline{\hspace{1.5cm}}
  \end{center}

  \item \textbf{Leia e escreva o número de sílabas de cada palavra:}
  \begin{center}
  CARRO $\rightarrow$ \underline{\hspace{0.9cm}} \quad TERRA $\rightarrow$ \underline{\hspace{0.9cm}} \quad BARRO $\rightarrow$ \underline{\hspace{0.9cm}} \quad SORRISO $\rightarrow$ \underline{\hspace{0.9cm}}
  \end{center}

  \item \textbf{Copie as frases em letra cursiva:}
  \begin{center}
  O carro andou na terra barrenta.\\
  \underline{\hspace{12cm}}\\
  O sorriso dela é bonito.\\
  \underline{\hspace{12cm}}
  \end{center}

  \item \textbf{Desenhe a frase:}
  \begin{center}
  \begin{tikzpicture}
    \draw[thick, dashed] (0,0) rectangle (12,4);
  \end{tikzpicture}
  \end{center}

\end{enumerate}

\end{exercicio}

%% COMPLEMENTAR C-34: REVISÃO S, Z E Ç

\plancheta{Folha Complementar C-34 --- Reforço da Folha E-34 (Revisão S, Z e Ç)}

\begin{exercicio}[Folha Complementar C-34 --- Reforço da Folha E-34 (Revisão S, Z e Ç)]

\metadata{C-34}{Silábico-Alfabético}{EF02LP03, EF02LP04, EF02LP07}{D1}{EF02LP04}{CNE/CEB nº 2/2012}

\kumontiming{2}{90}

\noindent\textbf{Nome:} \underline{\hspace{3.2cm}} \quad \textbf{Data:} \underline{\hspace{1.6cm}} \quad \textbf{Nota:} \underline{\hspace{0.9cm}}/10


\begin{enumerate}

  \item \textbf{Copie as palavras em letra cursiva:}
  \begin{center}
  casa \quad mesa \quad zero \quad zebra
  \end{center}
  \begin{center}
  \underline{\hspace{2.4cm}} \quad \underline{\hspace{2.4cm}} \quad \underline{\hspace{2.4cm}} \quad \underline{\hspace{2.4cm}}
  \end{center}

  \item \textbf{Escreva as sílabas das palavras:}
  \begin{center}
  CASA $=$ \underline{\hspace{1.5cm}} \underline{\hspace{1.5cm}} \quad MESA $=$ \underline{\hspace{1.5cm}} \underline{\hspace{1.5cm}} \quad ZERO $=$ \underline{\hspace{1.5cm}} \underline{\hspace{1.5cm}}
  \end{center}

  \item \textbf{Leia e escreva o número de sílabas de cada palavra:}
  \begin{center}
  CASA $\rightarrow$ \underline{\hspace{0.9cm}} \quad MESA $\rightarrow$ \underline{\hspace{0.9cm}} \quad ZERO $\rightarrow$ \underline{\hspace{0.9cm}} \quad ZEBRA $\rightarrow$ \underline{\hspace{0.9cm}}
  \end{center}

  \item \textbf{Copie as frases em letra cursiva:}
  \begin{center}
  A casa da zebra é perto da mesa.\\
  \underline{\hspace{12cm}}\\
  O zero é redondo como o laço.\\
  \underline{\hspace{12cm}}
  \end{center}

  \item \textbf{Desenhe a frase:}
  \begin{center}
  \begin{tikzpicture}
    \draw[thick, dashed] (0,0) rectangle (12,4);
  \end{tikzpicture}
  \end{center}

\end{enumerate}

\end{exercicio}

%% COMPLEMENTAR C-35: REVISÃO AM E AN

\plancheta{Folha Complementar C-35 --- Reforço da Folha E-35 (Revisão AM e AN)}

\begin{exercicio}[Folha Complementar C-35 --- Reforço da Folha E-35 (Revisão AM e AN)]

\metadata{C-35}{Silábico-Alfabético}{EF02LP03, EF02LP04, EF02LP07}{D1}{EF02LP04}{CNE/CEB nº 2/2012}

\kumontiming{2}{90}

\noindent\textbf{Nome:} \underline{\hspace{3.2cm}} \quad \textbf{Data:} \underline{\hspace{1.6cm}} \quad \textbf{Nota:} \underline{\hspace{0.9cm}}/10


\begin{enumerate}

  \item \textbf{Copie as palavras em letra cursiva:}
  \begin{center}
  também \quad sambar \quad canto \quad sangue
  \end{center}
  \begin{center}
  \underline{\hspace{2.4cm}} \quad \underline{\hspace{2.4cm}} \quad \underline{\hspace{2.4cm}} \quad \underline{\hspace{2.4cm}}
  \end{center}

  \item \textbf{Escreva as sílabas das palavras:}
  \begin{center}
  TAMBÉM $=$ \underline{\hspace{1.5cm}} \underline{\hspace{1.5cm}} \quad SAMBAR $=$ \underline{\hspace{1.5cm}} \underline{\hspace{1.5cm}} \quad CANTO $=$ \underline{\hspace{1.5cm}} \underline{\hspace{1.5cm}}
  \end{center}

  \item \textbf{Complete a tabela de sílabas:}

\begin{center}
\resizebox{\textwidth}{!}{%
\begin{tabular}{ccccc}
  \sil{AM}{A} & \sil{AM}{E} & \sil{AM}{I} & \sil{AM}{O} & \sil{AM}{U} \\
  \end{tabular}}
\end{center}

  \item \textbf{Copie as frases em letra cursiva:}
  \begin{center}
  Eu também quero sambar no canto.\\
  \underline{\hspace{12cm}}\\
  O banco fica no campo.\\
  \underline{\hspace{12cm}}
  \end{center}

  \item \textbf{Desenhe a frase:}
  \begin{center}
  \begin{tikzpicture}
    \draw[thick, dashed] (0,0) rectangle (12,4);
  \end{tikzpicture}
  \end{center}

\end{enumerate}

\end{exercicio}

%% COMPLEMENTAR C-36: REVISÃO X

\plancheta{Folha Complementar C-36 --- Reforço da Folha E-36 (Revisão X)}

\begin{exercicio}[Folha Complementar C-36 --- Reforço da Folha E-36 (Revisão X)]

\metadata{C-36}{Silábico-Alfabético}{EF02LP03, EF02LP04, EF02LP07}{D1}{EF02LP04}{CNE/CEB nº 2/2012}

\kumontiming{2}{90}

\noindent\textbf{Nome:} \underline{\hspace{3.2cm}} \quad \textbf{Data:} \underline{\hspace{1.6cm}} \quad \textbf{Nota:} \underline{\hspace{0.9cm}}/10


\begin{enumerate}

  \item \textbf{Copie as palavras em letra cursiva:}
  \begin{center}
  xadrez \quad peixe \quad caixa \quad exame
  \end{center}
  \begin{center}
  \underline{\hspace{2.4cm}} \quad \underline{\hspace{2.4cm}} \quad \underline{\hspace{2.4cm}} \quad \underline{\hspace{2.4cm}}
  \end{center}

  \item \textbf{Escreva as sílabas das palavras:}
  \begin{center}
  XADREZ $=$ \underline{\hspace{1.5cm}} \underline{\hspace{1.5cm}} \quad PEIXE $=$ \underline{\hspace{1.5cm}} \underline{\hspace{1.5cm}} \quad CAIXA $=$ \underline{\hspace{1.5cm}} \underline{\hspace{1.5cm}}
  \end{center}

  \item \textbf{Leia e escreva o número de sílabas de cada palavra:}
  \begin{center}
  XADREZ $\rightarrow$ \underline{\hspace{0.9cm}} \quad PEIXE $\rightarrow$ \underline{\hspace{0.9cm}} \quad CAIXA $\rightarrow$ \underline{\hspace{0.9cm}} \quad EXAME $\rightarrow$ \underline{\hspace{0.9cm}}
  \end{center}

  \item \textbf{Copie as frases em letra cursiva:}
  \begin{center}
  O peixe está na caixa do táxi.\\
  \underline{\hspace{12cm}}\\
  O xadrez é um jogo de atenção.\\
  \underline{\hspace{12cm}}
  \end{center}

  \item \textbf{Desenhe a frase:}
  \begin{center}
  \begin{tikzpicture}
    \draw[thick, dashed] (0,0) rectangle (12,4);
  \end{tikzpicture}
  \end{center}

\end{enumerate}

\end{exercicio}

%% COMPLEMENTAR C-37: LEITURA DE FRASES 1

\plancheta{Folha Complementar C-37 --- Reforço da Folha E-37 (Leitura de Frases 1)}

\begin{exercicio}[Folha Complementar C-37 --- Reforço da Folha E-37 (Leitura de Frases 1)]

\metadata{C-37}{Silábico-Alfabético}{EF02LP03, EF02LP04, EF02LP07}{D1}{EF02LP04}{CNE/CEB nº 2/2012}

\kumontiming{2}{90}

\noindent\textbf{Nome:} \underline{\hspace{3.2cm}} \quad \textbf{Data:} \underline{\hspace{1.6cm}} \quad \textbf{Nota:} \underline{\hspace{0.9cm}}/10


\begin{enumerate}

  \item \textbf{Copie as palavras em letra cursiva:}
  \begin{center}
  chuva \quad coelho \quad trem \quad flor
  \end{center}
  \begin{center}
  \underline{\hspace{2.4cm}} \quad \underline{\hspace{2.4cm}} \quad \underline{\hspace{2.4cm}} \quad \underline{\hspace{2.4cm}}
  \end{center}

  \item \textbf{Escreva as sílabas das palavras:}
  \begin{center}
  CHUVA $=$ \underline{\hspace{1.5cm}} \underline{\hspace{1.5cm}} \quad COELHO $=$ \underline{\hspace{1.5cm}} \underline{\hspace{1.5cm}} \quad TREM $=$ \underline{\hspace{1.5cm}} \underline{\hspace{1.5cm}}
  \end{center}

  \item \textbf{Leia e escreva o número de sílabas de cada palavra:}
  \begin{center}
  CHUVA $\rightarrow$ \underline{\hspace{0.9cm}} \quad COELHO $\rightarrow$ \underline{\hspace{0.9cm}} \quad TREM $\rightarrow$ \underline{\hspace{0.9cm}} \quad FLOR $\rightarrow$ \underline{\hspace{0.9cm}}
  \end{center}

  \item \textbf{Copie as frases em letra cursiva:}
  \begin{center}
  A chuva cai no telhado do coelho.\\
  \underline{\hspace{12cm}}\\
  O trem passa pela casa de flor.\\
  \underline{\hspace{12cm}}
  \end{center}

  \item \textbf{Desenhe a frase:}
  \begin{center}
  \begin{tikzpicture}
    \draw[thick, dashed] (0,0) rectangle (12,4);
  \end{tikzpicture}
  \end{center}

\end{enumerate}

\end{exercicio}

%% COMPLEMENTAR C-38: LEITURA DE FRASES 2

\plancheta{Folha Complementar C-38 --- Reforço da Folha E-38 (Leitura de Frases 2)}

\begin{exercicio}[Folha Complementar C-38 --- Reforço da Folha E-38 (Leitura de Frases 2)]

\metadata{C-38}{Silábico-Alfabético}{EF02LP03, EF02LP04, EF02LP07}{D1}{EF02LP04}{CNE/CEB nº 2/2012}

\kumontiming{2}{90}

\noindent\textbf{Nome:} \underline{\hspace{3.2cm}} \quad \textbf{Data:} \underline{\hspace{1.6cm}} \quad \textbf{Nota:} \underline{\hspace{0.9cm}}/10


\begin{enumerate}

  \item \textbf{Copie as palavras em letra cursiva:}
  \begin{center}
  biblioteca \quad história \quad capítulo \quad autor
  \end{center}
  \begin{center}
  \underline{\hspace{2.4cm}} \quad \underline{\hspace{2.4cm}} \quad \underline{\hspace{2.4cm}} \quad \underline{\hspace{2.4cm}}
  \end{center}

  \item \textbf{Escreva as sílabas das palavras:}
  \begin{center}
  BIBLIOTECA $=$ \underline{\hspace{1.5cm}} \underline{\hspace{1.5cm}} \quad HISTÓRIA $=$ \underline{\hspace{1.5cm}} \underline{\hspace{1.5cm}} \quad CAPÍTULO $=$ \underline{\hspace{1.5cm}} \underline{\hspace{1.5cm}}
  \end{center}

  \item \textbf{Leia e escreva o número de sílabas de cada palavra:}
  \begin{center}
  BIBLIOTECA $\rightarrow$ \underline{\hspace{0.9cm}} \quad HISTÓRIA $\rightarrow$ \underline{\hspace{0.9cm}} \quad CAPÍTULO $\rightarrow$ \underline{\hspace{0.9cm}} \quad AUTOR $\rightarrow$ \underline{\hspace{0.9cm}}
  \end{center}

  \item \textbf{Copie as frases em letra cursiva:}
  \begin{center}
  A biblioteca guarda a história.\\
  \underline{\hspace{12cm}}\\
  O capítulo novo começa na página 10.\\
  \underline{\hspace{12cm}}
  \end{center}

  \item \textbf{Desenhe a frase:}
  \begin{center}
  \begin{tikzpicture}
    \draw[thick, dashed] (0,0) rectangle (12,4);
  \end{tikzpicture}
  \end{center}

\end{enumerate}

\end{exercicio}

%% COMPLEMENTAR C-39: DITADO DAS GRAFIAS DO ANO

\plancheta{Folha Complementar C-39 --- Reforço da Folha E-39 (Ditado das Grafias do Ano)}

\begin{exercicio}[Folha Complementar C-39 --- Reforço da Folha E-39 (Ditado das Grafias do Ano)]

\metadata{C-39}{Silábico-Alfabético}{EF02LP03, EF02LP04, EF02LP07}{D1}{EF02LP04}{CNE/CEB nº 2/2012}

\kumontiming{2}{90}

\noindent\textbf{Nome:} \underline{\hspace{3.2cm}} \quad \textbf{Data:} \underline{\hspace{1.6cm}} \quad \textbf{Nota:} \underline{\hspace{0.9cm}}/10


\begin{enumerate}

  \item \textbf{Copie as palavras em letra cursiva:}
  \begin{center}
  chave \quad coelho \quad ninho \quad queijo
  \end{center}
  \begin{center}
  \underline{\hspace{2.4cm}} \quad \underline{\hspace{2.4cm}} \quad \underline{\hspace{2.4cm}} \quad \underline{\hspace{2.4cm}}
  \end{center}

  \item \textbf{Escreva as sílabas das palavras:}
  \begin{center}
  CHAVE $=$ \underline{\hspace{1.5cm}} \underline{\hspace{1.5cm}} \quad COELHO $=$ \underline{\hspace{1.5cm}} \underline{\hspace{1.5cm}} \quad NINHO $=$ \underline{\hspace{1.5cm}} \underline{\hspace{1.5cm}}
  \end{center}

  \item \textbf{Leia e escreva o número de sílabas de cada palavra:}
  \begin{center}
  CHAVE $\rightarrow$ \underline{\hspace{0.9cm}} \quad COELHO $\rightarrow$ \underline{\hspace{0.9cm}} \quad NINHO $\rightarrow$ \underline{\hspace{0.9cm}} \quad QUEIJO $\rightarrow$ \underline{\hspace{0.9cm}}
  \end{center}

  \item \textbf{Copie as frases em letra cursiva:}
  \begin{center}
  A chave abre o cofre do coelho.\\
  \underline{\hspace{12cm}}\\
  O ninho do queijo é de palha.\\
  \underline{\hspace{12cm}}
  \end{center}

  \item \textbf{Desenhe a frase:}
  \begin{center}
  \begin{tikzpicture}
    \draw[thick, dashed] (0,0) rectangle (12,4);
  \end{tikzpicture}
  \end{center}

\end{enumerate}

\end{exercicio}

%% COMPLEMENTAR C-40: PROVA DE FLUÊNCIA LEITORA

\plancheta{Folha Complementar C-40 --- Reforço da Folha E-40 (Prova de Fluência Leitora)}

\begin{exercicio}[Folha Complementar C-40 --- Reforço da Folha E-40 (Prova de Fluência Leitora)]

\metadata{C-40}{Silábico-Alfabético}{EF02LP03, EF02LP04, EF02LP07}{D1}{EF02LP04}{CNE/CEB nº 2/2012}

\kumontiming{2}{90}

\noindent\textbf{Nome:} \underline{\hspace{3.2cm}} \quad \textbf{Data:} \underline{\hspace{1.6cm}} \quad \textbf{Nota:} \underline{\hspace{0.9cm}}/10


\begin{enumerate}

  \item \textbf{Copie as palavras em letra cursiva:}
  \begin{center}
  leitura \quad fluência \quad entonação \quad precisão
  \end{center}
  \begin{center}
  \underline{\hspace{2.4cm}} \quad \underline{\hspace{2.4cm}} \quad \underline{\hspace{2.4cm}} \quad \underline{\hspace{2.4cm}}
  \end{center}

  \item \textbf{Escreva as sílabas das palavras:}
  \begin{center}
  LEITURA $=$ \underline{\hspace{1.5cm}} \underline{\hspace{1.5cm}} \quad FLUÊNCIA $=$ \underline{\hspace{1.5cm}} \underline{\hspace{1.5cm}} \quad ENTONAÇÃO $=$ \underline{\hspace{1.5cm}} \underline{\hspace{1.5cm}}
  \end{center}

  \item \textbf{Leia e escreva o número de sílabas de cada palavra:}
  \begin{center}
  LEITURA $\rightarrow$ \underline{\hspace{0.9cm}} \quad FLUÊNCIA $\rightarrow$ \underline{\hspace{0.9cm}} \quad ENTONAÇÃO $\rightarrow$ \underline{\hspace{0.9cm}} \quad PRECISÃO $\rightarrow$ \underline{\hspace{0.9cm}}
  \end{center}

  \item \textbf{Copie as frases em letra cursiva:}
  \begin{center}
  A leitura com fluência tem ritmo e entonação.\\
  \underline{\hspace{12cm}}\\
  A precisão vem do treino diário.\\
  \underline{\hspace{12cm}}
  \end{center}

  \item \textbf{Desenhe a frase:}
  \begin{center}
  \begin{tikzpicture}
    \draw[thick, dashed] (0,0) rectangle (12,4);
  \end{tikzpicture}
  \end{center}

\end{enumerate}

\end{exercicio}

%% COMPLEMENTAR C-41: CALIGRAFIA CURSIVA 1

\plancheta{Folha Complementar C-41 --- Reforço da Folha E-41 (Caligrafia Cursiva 1)}

\begin{exercicio}[Folha Complementar C-41 --- Reforço da Folha E-41 (Caligrafia Cursiva 1)]

\metadata{C-41}{Silábico-Alfabético}{EF02LP03, EF02LP04, EF02LP07}{D1}{EF02LP04}{CNE/CEB nº 2/2012}

\kumontiming{2}{90}

\noindent\textbf{Nome:} \underline{\hspace{3.2cm}} \quad \textbf{Data:} \underline{\hspace{1.6cm}} \quad \textbf{Nota:} \underline{\hspace{0.9cm}}/10


\begin{enumerate}

  \item \textbf{Copie as palavras em letra cursiva:}
  \begin{center}
  chave \quad coelho \quad ninho \quad queijo
  \end{center}
  \begin{center}
  \underline{\hspace{2.4cm}} \quad \underline{\hspace{2.4cm}} \quad \underline{\hspace{2.4cm}} \quad \underline{\hspace{2.4cm}}
  \end{center}

  \item \textbf{Escreva as sílabas das palavras:}
  \begin{center}
  CHAVE $=$ \underline{\hspace{1.5cm}} \underline{\hspace{1.5cm}} \quad COELHO $=$ \underline{\hspace{1.5cm}} \underline{\hspace{1.5cm}} \quad NINHO $=$ \underline{\hspace{1.5cm}} \underline{\hspace{1.5cm}}
  \end{center}

  \item \textbf{Leia e escreva o número de sílabas de cada palavra:}
  \begin{center}
  CHAVE $\rightarrow$ \underline{\hspace{0.9cm}} \quad COELHO $\rightarrow$ \underline{\hspace{0.9cm}} \quad NINHO $\rightarrow$ \underline{\hspace{0.9cm}} \quad QUEIJO $\rightarrow$ \underline{\hspace{0.9cm}}
  \end{center}

  \item \textbf{Copie as frases em letra cursiva:}
  \begin{center}
  A chave do coelho abriu o ninho.\\
  \underline{\hspace{12cm}}\\
  A guerra do carro passou na casa.\\
  \underline{\hspace{12cm}}
  \end{center}

  \item \textbf{Desenhe a frase:}
  \begin{center}
  \begin{tikzpicture}
    \draw[thick, dashed] (0,0) rectangle (12,4);
  \end{tikzpicture}
  \end{center}

\end{enumerate}

\end{exercicio}

%% COMPLEMENTAR C-42: CALIGRAFIA CURSIVA 2

\plancheta{Folha Complementar C-42 --- Reforço da Folha E-42 (Caligrafia Cursiva 2)}

\begin{exercicio}[Folha Complementar C-42 --- Reforço da Folha E-42 (Caligrafia Cursiva 2)]

\metadata{C-42}{Silábico-Alfabético}{EF02LP03, EF02LP04, EF02LP07}{D1}{EF02LP04}{CNE/CEB nº 2/2012}

\kumontiming{2}{90}

\noindent\textbf{Nome:} \underline{\hspace{3.2cm}} \quad \textbf{Data:} \underline{\hspace{1.6cm}} \quad \textbf{Nota:} \underline{\hspace{0.9cm}}/10


\begin{enumerate}

  \item \textbf{Copie as palavras em letra cursiva:}
  \begin{center}
  biblioteca \quad estante \quad revista \quad história
  \end{center}
  \begin{center}
  \underline{\hspace{2.4cm}} \quad \underline{\hspace{2.4cm}} \quad \underline{\hspace{2.4cm}} \quad \underline{\hspace{2.4cm}}
  \end{center}

  \item \textbf{Escreva as sílabas das palavras:}
  \begin{center}
  BIBLIOTECA $=$ \underline{\hspace{1.5cm}} \underline{\hspace{1.5cm}} \quad ESTANTE $=$ \underline{\hspace{1.5cm}} \underline{\hspace{1.5cm}} \quad REVISTA $=$ \underline{\hspace{1.5cm}} \underline{\hspace{1.5cm}}
  \end{center}

  \item \textbf{Leia e escreva o número de sílabas de cada palavra:}
  \begin{center}
  BIBLIOTECA $\rightarrow$ \underline{\hspace{0.9cm}} \quad ESTANTE $\rightarrow$ \underline{\hspace{0.9cm}} \quad REVISTA $\rightarrow$ \underline{\hspace{0.9cm}} \quad HISTÓRIA $\rightarrow$ \underline{\hspace{0.9cm}}
  \end{center}

  \item \textbf{Copie as frases em letra cursiva:}
  \begin{center}
  A biblioteca da escola tem estante nova.\\
  \underline{\hspace{12cm}}\\
  A história do capítulo é de um autor famoso.\\
  \underline{\hspace{12cm}}
  \end{center}

  \item \textbf{Desenhe a frase:}
  \begin{center}
  \begin{tikzpicture}
    \draw[thick, dashed] (0,0) rectangle (12,4);
  \end{tikzpicture}
  \end{center}

\end{enumerate}

\end{exercicio}

%% COMPLEMENTAR C-43: AUTODITADO

\plancheta{Folha Complementar C-43 --- Reforço da Folha E-43 (Autoditado)}

\begin{exercicio}[Folha Complementar C-43 --- Reforço da Folha E-43 (Autoditado)]

\metadata{C-43}{Silábico-Alfabético}{EF02LP03, EF02LP04, EF02LP07}{D1}{EF02LP04}{CNE/CEB nº 2/2012}

\kumontiming{2}{90}

\noindent\textbf{Nome:} \underline{\hspace{3.2cm}} \quad \textbf{Data:} \underline{\hspace{1.6cm}} \quad \textbf{Nota:} \underline{\hspace{0.9cm}}/10


\begin{enumerate}

  \item \textbf{Copie as palavras em letra cursiva:}
  \begin{center}
  chave \quad coelho \quad ninho \quad queijo
  \end{center}
  \begin{center}
  \underline{\hspace{2.4cm}} \quad \underline{\hspace{2.4cm}} \quad \underline{\hspace{2.4cm}} \quad \underline{\hspace{2.4cm}}
  \end{center}

  \item \textbf{Escreva as sílabas das palavras:}
  \begin{center}
  CHAVE $=$ \underline{\hspace{1.5cm}} \underline{\hspace{1.5cm}} \quad COELHO $=$ \underline{\hspace{1.5cm}} \underline{\hspace{1.5cm}} \quad NINHO $=$ \underline{\hspace{1.5cm}} \underline{\hspace{1.5cm}}
  \end{center}

  \item \textbf{Leia e escreva o número de sílabas de cada palavra:}
  \begin{center}
  CHAVE $\rightarrow$ \underline{\hspace{0.9cm}} \quad COELHO $\rightarrow$ \underline{\hspace{0.9cm}} \quad NINHO $\rightarrow$ \underline{\hspace{0.9cm}} \quad QUEIJO $\rightarrow$ \underline{\hspace{0.9cm}}
  \end{center}

  \item \textbf{Copie as frases em letra cursiva:}
  \begin{center}
  A chave do coelho é dourada.\\
  \underline{\hspace{12cm}}\\
  O ninho do queijo é de palha.\\
  \underline{\hspace{12cm}}
  \end{center}

  \item \textbf{Desenhe a frase:}
  \begin{center}
  \begin{tikzpicture}
    \draw[thick, dashed] (0,0) rectangle (12,4);
  \end{tikzpicture}
  \end{center}

\end{enumerate}

\end{exercicio}

%% COMPLEMENTAR C-44: BANCO DE PALAVRAS

\plancheta{Folha Complementar C-44 --- Reforço da Folha E-44 (Banco de Palavras)}

\begin{exercicio}[Folha Complementar C-44 --- Reforço da Folha E-44 (Banco de Palavras)]

\metadata{C-44}{Silábico-Alfabético}{EF02LP03, EF02LP04, EF02LP07}{D1}{EF02LP04}{CNE/CEB nº 2/2012}

\kumontiming{2}{90}

\noindent\textbf{Nome:} \underline{\hspace{3.2cm}} \quad \textbf{Data:} \underline{\hspace{1.6cm}} \quad \textbf{Nota:} \underline{\hspace{0.9cm}}/10


\begin{enumerate}

  \item \textbf{Copie as palavras em letra cursiva:}
  \begin{center}
  chuva \quad chapéu \quad ilha \quad orelha
  \end{center}
  \begin{center}
  \underline{\hspace{2.4cm}} \quad \underline{\hspace{2.4cm}} \quad \underline{\hspace{2.4cm}} \quad \underline{\hspace{2.4cm}}
  \end{center}

  \item \textbf{Escreva as sílabas das palavras:}
  \begin{center}
  CHUVA $=$ \underline{\hspace{1.5cm}} \underline{\hspace{1.5cm}} \quad CHAPÉU $=$ \underline{\hspace{1.5cm}} \underline{\hspace{1.5cm}} \quad ILHA $=$ \underline{\hspace{1.5cm}} \underline{\hspace{1.5cm}}
  \end{center}

  \item \textbf{Leia e escreva o número de sílabas de cada palavra:}
  \begin{center}
  CHUVA $\rightarrow$ \underline{\hspace{0.9cm}} \quad CHAPÉU $\rightarrow$ \underline{\hspace{0.9cm}} \quad ILHA $\rightarrow$ \underline{\hspace{0.9cm}} \quad ORELHA $\rightarrow$ \underline{\hspace{0.9cm}}
  \end{center}

  \item \textbf{Copie as frases em letra cursiva:}
  \begin{center}
  A chuva molhou o chapéu.\\
  \underline{\hspace{12cm}}\\
  A ilha tem a orelha do mapa.\\
  \underline{\hspace{12cm}}
  \end{center}

  \item \textbf{Desenhe a frase:}
  \begin{center}
  \begin{tikzpicture}
    \draw[thick, dashed] (0,0) rectangle (12,4);
  \end{tikzpicture}
  \end{center}

\end{enumerate}

\end{exercicio}

%% COMPLEMENTAR C-45: LISTAS E COMPLETAR FRASES

\plancheta{Folha Complementar C-45 --- Reforço da Folha E-45 (Listas e Completar Frases)}

\begin{exercicio}[Folha Complementar C-45 --- Reforço da Folha E-45 (Listas e Completar Frases)]

\metadata{C-45}{Silábico-Alfabético}{EF02LP03, EF02LP04, EF02LP07}{D1}{EF02LP04}{CNE/CEB nº 2/2012}

\kumontiming{2}{90}

\noindent\textbf{Nome:} \underline{\hspace{3.2cm}} \quad \textbf{Data:} \underline{\hspace{1.6cm}} \quad \textbf{Nota:} \underline{\hspace{0.9cm}}/10


\begin{enumerate}

  \item \textbf{Copie as palavras em letra cursiva:}
  \begin{center}
  braço \quad trem \quad fruta \quad flor
  \end{center}
  \begin{center}
  \underline{\hspace{2.4cm}} \quad \underline{\hspace{2.4cm}} \quad \underline{\hspace{2.4cm}} \quad \underline{\hspace{2.4cm}}
  \end{center}

  \item \textbf{Escreva as sílabas das palavras:}
  \begin{center}
  BRAÇO $=$ \underline{\hspace{1.5cm}} \underline{\hspace{1.5cm}} \quad TREM $=$ \underline{\hspace{1.5cm}} \underline{\hspace{1.5cm}} \quad FRUTA $=$ \underline{\hspace{1.5cm}} \underline{\hspace{1.5cm}}
  \end{center}

  \item \textbf{Leia e escreva o número de sílabas de cada palavra:}
  \begin{center}
  BRAÇO $\rightarrow$ \underline{\hspace{0.9cm}} \quad TREM $\rightarrow$ \underline{\hspace{0.9cm}} \quad FRUTA $\rightarrow$ \underline{\hspace{0.9cm}} \quad FLOR $\rightarrow$ \underline{\hspace{0.9cm}}
  \end{center}

  \item \textbf{Copie as frases em letra cursiva:}
  \begin{center}
  O braço aponta para o trem.\\
  \underline{\hspace{12cm}}\\
  A fruta e a flor enfeitam o prato.\\
  \underline{\hspace{12cm}}
  \end{center}

  \item \textbf{Desenhe a frase:}
  \begin{center}
  \begin{tikzpicture}
    \draw[thick, dashed] (0,0) rectangle (12,4);
  \end{tikzpicture}
  \end{center}

\end{enumerate}

\end{exercicio}

%% COMPLEMENTAR C-46: COMPLETAR FRASES

\plancheta{Folha Complementar C-46 --- Reforço da Folha E-46 (Completar Frases)}

\begin{exercicio}[Folha Complementar C-46 --- Reforço da Folha E-46 (Completar Frases)]

\metadata{C-46}{Silábico-Alfabético}{EF02LP03, EF02LP04, EF02LP07}{D1}{EF02LP04}{CNE/CEB nº 2/2012}

\kumontiming{2}{90}

\noindent\textbf{Nome:} \underline{\hspace{3.2cm}} \quad \textbf{Data:} \underline{\hspace{1.6cm}} \quad \textbf{Nota:} \underline{\hspace{0.9cm}}/10


\begin{enumerate}

  \item \textbf{Copie as palavras em letra cursiva:}
  \begin{center}
  professora \quad escola \quad amigo \quad pátio
  \end{center}
  \begin{center}
  \underline{\hspace{2.4cm}} \quad \underline{\hspace{2.4cm}} \quad \underline{\hspace{2.4cm}} \quad \underline{\hspace{2.4cm}}
  \end{center}

  \item \textbf{Escreva as sílabas das palavras:}
  \begin{center}
  PROFESSORA $=$ \underline{\hspace{1.5cm}} \underline{\hspace{1.5cm}} \quad ESCOLA $=$ \underline{\hspace{1.5cm}} \underline{\hspace{1.5cm}} \quad AMIGO $=$ \underline{\hspace{1.5cm}} \underline{\hspace{1.5cm}}
  \end{center}

  \item \textbf{Leia e escreva o número de sílabas de cada palavra:}
  \begin{center}
  PROFESSORA $\rightarrow$ \underline{\hspace{0.9cm}} \quad ESCOLA $\rightarrow$ \underline{\hspace{0.9cm}} \quad AMIGO $\rightarrow$ \underline{\hspace{0.9cm}} \quad PÁTIO $\rightarrow$ \underline{\hspace{0.9cm}}
  \end{center}

  \item \textbf{Copie as frases em letra cursiva:}
  \begin{center}
  A professora escreve no quadro da escola.\\
  \underline{\hspace{12cm}}\\
  O amigo lancha no pátio.\\
  \underline{\hspace{12cm}}
  \end{center}

  \item \textbf{Desenhe a frase:}
  \begin{center}
  \begin{tikzpicture}
    \draw[thick, dashed] (0,0) rectangle (12,4);
  \end{tikzpicture}
  \end{center}

\end{enumerate}

\end{exercicio}

%% COMPLEMENTAR C-47: LEITURA DE TEXTO CURTO 1

\plancheta{Folha Complementar C-47 --- Reforço da Folha E-47 (Leitura de Texto Curto 1)}

\begin{exercicio}[Folha Complementar C-47 --- Reforço da Folha E-47 (Leitura de Texto Curto 1)]

\metadata{C-47}{Silábico-Alfabético}{EF02LP03, EF02LP04, EF02LP07}{D1}{EF02LP04}{CNE/CEB nº 2/2012}

\kumontiming{2}{90}

\noindent\textbf{Nome:} \underline{\hspace{3.2cm}} \quad \textbf{Data:} \underline{\hspace{1.6cm}} \quad \textbf{Nota:} \underline{\hspace{0.9cm}}/10


\begin{enumerate}

  \item \textbf{Copie as palavras em letra cursiva:}
  \begin{center}
  chuva \quad telhado \quad gato \quad janela
  \end{center}
  \begin{center}
  \underline{\hspace{2.4cm}} \quad \underline{\hspace{2.4cm}} \quad \underline{\hspace{2.4cm}} \quad \underline{\hspace{2.4cm}}
  \end{center}

  \item \textbf{Escreva as sílabas das palavras:}
  \begin{center}
  CHUVA $=$ \underline{\hspace{1.5cm}} \underline{\hspace{1.5cm}} \quad TELHADO $=$ \underline{\hspace{1.5cm}} \underline{\hspace{1.5cm}} \quad GATO $=$ \underline{\hspace{1.5cm}} \underline{\hspace{1.5cm}}
  \end{center}

  \item \textbf{Leia e escreva o número de sílabas de cada palavra:}
  \begin{center}
  CHUVA $\rightarrow$ \underline{\hspace{0.9cm}} \quad TELHADO $\rightarrow$ \underline{\hspace{0.9cm}} \quad GATO $\rightarrow$ \underline{\hspace{0.9cm}} \quad JANELA $\rightarrow$ \underline{\hspace{0.9cm}}
  \end{center}

  \item \textbf{Copie as frases em letra cursiva:}
  \begin{center}
  A chuva cai no telhado.\\
  \underline{\hspace{12cm}}\\
  O gato olha pela janela.\\
  \underline{\hspace{12cm}}
  \end{center}

  \item \textbf{Desenhe a frase:}
  \begin{center}
  \begin{tikzpicture}
    \draw[thick, dashed] (0,0) rectangle (12,4);
  \end{tikzpicture}
  \end{center}

\end{enumerate}

\end{exercicio}

%% COMPLEMENTAR C-48: LEITURA DE TEXTO CURTO 2

\plancheta{Folha Complementar C-48 --- Reforço da Folha E-48 (Leitura de Texto Curto 2)}

\begin{exercicio}[Folha Complementar C-48 --- Reforço da Folha E-48 (Leitura de Texto Curto 2)]

\metadata{C-48}{Silábico-Alfabético}{EF02LP03, EF02LP04, EF02LP07}{D1}{EF02LP04}{CNE/CEB nº 2/2012}

\kumontiming{2}{90}

\noindent\textbf{Nome:} \underline{\hspace{3.2cm}} \quad \textbf{Data:} \underline{\hspace{1.6cm}} \quad \textbf{Nota:} \underline{\hspace{0.9cm}}/10


\begin{enumerate}

  \item \textbf{Copie as palavras em letra cursiva:}
  \begin{center}
  menino \quad pipas \quad vento \quad campo
  \end{center}
  \begin{center}
  \underline{\hspace{2.4cm}} \quad \underline{\hspace{2.4cm}} \quad \underline{\hspace{2.4cm}} \quad \underline{\hspace{2.4cm}}
  \end{center}

  \item \textbf{Escreva as sílabas das palavras:}
  \begin{center}
  MENINO $=$ \underline{\hspace{1.5cm}} \underline{\hspace{1.5cm}} \quad PIPAS $=$ \underline{\hspace{1.5cm}} \underline{\hspace{1.5cm}} \quad VENTO $=$ \underline{\hspace{1.5cm}} \underline{\hspace{1.5cm}}
  \end{center}

  \item \textbf{Leia e escreva o número de sílabas de cada palavra:}
  \begin{center}
  MENINO $\rightarrow$ \underline{\hspace{0.9cm}} \quad PIPAS $\rightarrow$ \underline{\hspace{0.9cm}} \quad VENTO $\rightarrow$ \underline{\hspace{0.9cm}} \quad CAMPO $\rightarrow$ \underline{\hspace{0.9cm}}
  \end{center}

  \item \textbf{Copie as frases em letra cursiva:}
  \begin{center}
  O menino solta pipas no campo.\\
  \underline{\hspace{12cm}}\\
  O vento leva a nuvem pelo céu.\\
  \underline{\hspace{12cm}}
  \end{center}

  \item \textbf{Desenhe a frase:}
  \begin{center}
  \begin{tikzpicture}
    \draw[thick, dashed] (0,0) rectangle (12,4);
  \end{tikzpicture}
  \end{center}

\end{enumerate}

\end{exercicio}

%% COMPLEMENTAR C-49: PRODUÇÃO GUIADA

\plancheta{Folha Complementar C-49 --- Reforço da Folha E-49 (Produção Guiada)}

\begin{exercicio}[Folha Complementar C-49 --- Reforço da Folha E-49 (Produção Guiada)]

\metadata{C-49}{Silábico-Alfabético}{EF02LP03, EF02LP04, EF02LP07}{D1}{EF02LP04}{CNE/CEB nº 2/2012}

\kumontiming{2}{90}

\noindent\textbf{Nome:} \underline{\hspace{3.2cm}} \quad \textbf{Data:} \underline{\hspace{1.6cm}} \quad \textbf{Nota:} \underline{\hspace{0.9cm}}/10


\begin{enumerate}

  \item \textbf{Copie as palavras em letra cursiva:}
  \begin{center}
  conto \quad fábula \quad bilhete \quad notícia
  \end{center}
  \begin{center}
  \underline{\hspace{2.4cm}} \quad \underline{\hspace{2.4cm}} \quad \underline{\hspace{2.4cm}} \quad \underline{\hspace{2.4cm}}
  \end{center}

  \item \textbf{Escreva as sílabas das palavras:}
  \begin{center}
  CONTO $=$ \underline{\hspace{1.5cm}} \underline{\hspace{1.5cm}} \quad FÁBULA $=$ \underline{\hspace{1.5cm}} \underline{\hspace{1.5cm}} \quad BILHETE $=$ \underline{\hspace{1.5cm}} \underline{\hspace{1.5cm}}
  \end{center}

  \item \textbf{Leia e escreva o número de sílabas de cada palavra:}
  \begin{center}
  CONTO $\rightarrow$ \underline{\hspace{0.9cm}} \quad FÁBULA $\rightarrow$ \underline{\hspace{0.9cm}} \quad BILHETE $\rightarrow$ \underline{\hspace{0.9cm}} \quad NOTÍCIA $\rightarrow$ \underline{\hspace{0.9cm}}
  \end{center}

  \item \textbf{Copie as frases em letra cursiva:}
  \begin{center}
  O conto virou fábula com moral.\\
  \underline{\hspace{12cm}}\\
  O bilhete virou notícia no jornal.\\
  \underline{\hspace{12cm}}
  \end{center}

  \item \textbf{Desenhe a frase:}
  \begin{center}
  \begin{tikzpicture}
    \draw[thick, dashed] (0,0) rectangle (12,4);
  \end{tikzpicture}
  \end{center}

\end{enumerate}

\end{exercicio}

%% COMPLEMENTAR C-50: REVISÃO GERAL FINAL

\plancheta{Folha Complementar C-50 --- Reforço da Folha E-50 (Revisão Geral Final)}

\begin{exercicio}[Folha Complementar C-50 --- Reforço da Folha E-50 (Revisão Geral Final)]

\metadata{C-50}{Silábico-Alfabético}{EF02LP03, EF02LP04, EF02LP07}{D1}{EF02LP04}{CNE/CEB nº 2/2012}

\kumontiming{2}{90}

\noindent\textbf{Nome:} \underline{\hspace{3.2cm}} \quad \textbf{Data:} \underline{\hspace{1.6cm}} \quad \textbf{Nota:} \underline{\hspace{0.9cm}}/10


\begin{enumerate}

  \item \textbf{Copie as palavras em letra cursiva:}
  \begin{center}
  chave \quad coelho \quad carro \quad casa
  \end{center}
  \begin{center}
  \underline{\hspace{2.4cm}} \quad \underline{\hspace{2.4cm}} \quad \underline{\hspace{2.4cm}} \quad \underline{\hspace{2.4cm}}
  \end{center}

  \item \textbf{Escreva as sílabas das palavras:}
  \begin{center}
  CHAVE $=$ \underline{\hspace{1.5cm}} \underline{\hspace{1.5cm}} \quad COELHO $=$ \underline{\hspace{1.5cm}} \underline{\hspace{1.5cm}} \quad CARRO $=$ \underline{\hspace{1.5cm}} \underline{\hspace{1.5cm}}
  \end{center}

  \item \textbf{Leia e escreva o número de sílabas de cada palavra:}
  \begin{center}
  CHAVE $\rightarrow$ \underline{\hspace{0.9cm}} \quad COELHO $\rightarrow$ \underline{\hspace{0.9cm}} \quad CARRO $\rightarrow$ \underline{\hspace{0.9cm}} \quad CASA $\rightarrow$ \underline{\hspace{0.9cm}}
  \end{center}

  \item \textbf{Copie as frases em letra cursiva:}
  \begin{center}
  A chave abriu o cofre do coelho.\\
  \underline{\hspace{12cm}}\\
  O carro levou a casa ao sarau.\\
  \underline{\hspace{12cm}}
  \end{center}

  \item \textbf{Desenhe a frase:}
  \begin{center}
  \begin{tikzpicture}
    \draw[thick, dashed] (0,0) rectangle (12,4);
  \end{tikzpicture}
  \end{center}

\end{enumerate}

\end{exercicio}
